% ============================================================================
%  INFORME FINAL — Migración MySQL 5.7 -> PostgreSQL 16
%  Cooperativa de Ahorro y Crédito Horizonte PyME S.R.L.
%  ----------------------------------------------------------------------------
%  Regla de oro del proyecto:
%    - formato/datos-institucion.sty -> datos compartidos (institución, logo)
%    - formato/preambulo.sty         -> toda la configuración (APA + Arial)
%    - formato/portada.tex           -> portada (datos institucionales)
%    - formato/referencias.bib       -> bibliografía BibTeX
%    - contenido/*.tex               -> el contenido real del informe
%    - anexos/*.tex                  -> material complementario
%    - main.tex                      -> integración + datos del documento
%
%  Compilar (desde la raíz de docs/):
%    make
% ============================================================================

\documentclass[12pt,a4paper]{article}

\usepackage{formato/preambulo}
\usepackage{formato/datos-institucion}

% ============================================================================
%  DATOS DEL DOCUMENTO
% ============================================================================
\newcommand{\tituloDocumento}{Migración heterogénea y endurecimiento de una base de datos financiera: de MySQL 5.7 a PostgreSQL 16 con Docker Compose}
\newcommand{\equipo}{Sudoers}
\newcommand{\integrantes}{%
  Joaquín Alessandro Felipez Rojas \par
  David Willy Cruz Huanca \par
  Nicolás Sebastián Reguerín Meneses \par
  Gustavo Abel Ignacio Mejía}
\newcommand{\asignatura}{Tecnología de Base de Datos I}
\newcommand{\docente}{Lineth Cindy Coca Rojas}
\newcommand{\fechaEntrega}{20/11/2026}

\begin{document}

\pagenumbering{roman}
% ============================================================================
%  PORTADA  (formato APA 7.ª edición)
%  ----------------------------------------------------------------------------
%  Los datos provienen de las macros definidas en main.tex. Editar ahí.
% ============================================================================

\begin{titlepage}
  \thispagestyle{empty}
  \setstretch{1}
  \setlength{\parskip}{0pt}

  \begin{center}
    \vspace*{0.6cm}

    {\bfseries\fontsize{20}{24}\selectfont \universidad \par}
    {\bfseries\fontsize{18}{22}\selectfont \facultad \par}
    {\bfseries\fontsize{16}{19}\selectfont \carrera \par}

    \vspace{0.8cm}

    % Logo de la institución (colocar "logo-upds.png" u otro nombre en figuras/)
    \IfFileExists{figuras/logo.png}{%
      \includegraphics[width=0.30\textwidth]{figuras/logo.png}\par
    }{}

    \vspace{0.9cm}

    {\bfseries\Large \tituloDocumento \par}

    \vspace{1.0cm}

    {\large \integrantes \par}

    \vspace{1.0cm}

    {\normalsize \textbf{Asignatura:} \asignatura \par}
    {\normalsize \textbf{Docente:} \docente \par}
    {\normalsize \textbf{Fecha de entrega:} \fechaEntrega \par}

    \vspace{1.1cm}

    {\normalsize \textbf{Equipo:} \equipo \par}

    \vspace{0.9cm}

    {\normalsize \ciudad \par}

    \vspace{0.2cm}

    {\normalsize \the\year \par}
  \end{center}
\end{titlepage}

\newpage
{\indicesCompactos\tableofcontents}
\newpage
{\indicesCompactos\listoftables}
\newpage
{\indicesCompactos\listoffigures}
\clearpage

\pagenumbering{arabic}

% --- CAPÍTULO I --------------------------------------------------------------
\caratulaCapitulo{CAPÍTULO I}{INTRODUCCIÓN Y CLIENTE}
% ============================================================================
%  CAPÍTULO I: INTRODUCCIÓN, CLIENTE Y ALCANCE
% ============================================================================

\section{Resumen ejecutivo}

La Cooperativa de Ahorro y Crédito Horizonte PyME S.R.L. gestiona sus créditos
y ahorros sobre una base de datos MySQL 5.7, motor que perdió el soporte
oficial en octubre de 2023. Para resolver ese riesgo se migró la base a
PostgreSQL 16, en un entorno desplegado íntegramente sobre contenedores Docker
y orquestado con Docker Compose. La migración incluyó el diagnóstico del
origen, el rediseño del esquema con tipos nativos, un proceso de carga
reproducible, el endurecimiento (\emph{hardening}) del servidor destino, roles
de mínimo privilegio, enmascaramiento de datos personales, respaldo cifrado y
pruebas automatizadas de cada resultado.

La migración se ejecutó sobre un conjunto de datos realista y determinista:
8 tablas, 2 vistas, 19 entradas de índice, 55 socios, 66 cuentas de ahorro,
50 solicitudes de crédito, 32 préstamos, 516 cuotas de amortización y 740
transacciones, para un total de 1\,469 filas migradas. El origen ocupa
0.28 MB y presenta incompatibilidades deliberadas con PostgreSQL: 7 columnas
\texttt{ENUM}, 10 columnas \texttt{DATETIME}, 4 columnas \texttt{TINYINT(1)},
claves \texttt{INT UNSIGNED AUTO\_INCREMENT}, collation \texttt{utf8\_general\_ci}
y ausencia total de claves foráneas. Además, el usuario de aplicación
\texttt{app\_legacy\_user} conserva un \texttt{GRANT ALL} sobre toda la base,
lo que constituye un hallazgo de auditoría de primer orden.

El destino se diseñó con \texttt{BOOLEAN}, \texttt{TIMESTAMPTZ},
identidades \texttt{BIGINT GENERATED ALWAYS AS IDENTITY}, restricciones
\texttt{CHECK}, 9 claves foráneas y 10 índices orientados a los reportes de
cartera y saldos. El endurecimiento incluye TLS/SSL obligatorio con
\texttt{hostssl}, autenticación \texttt{scram-sha-256}, revocación de
privilegios de \texttt{PUBLIC}, tres roles granulares (\texttt{admin\_dba},
\texttt{app\_core\_rw} y \texttt{auditor\_ro}), una vista enmascarada
\texttt{vw\_socios\_audit} y registro de conexiones, desconexiones y DDL.
El respaldo del destino se genera con \texttt{pg\_dump -Fc} y se cifra con
AES-256-CBC y PBKDF2; la restauración se probó y devolvió los 55 socios
esperados.

La verificación final arrojó 75 pruebas automatizadas en estado \texttt{PASS},
ninguna en \texttt{FAIL}, cero huérfanos en las 9 relaciones, conteos idénticos
entre origen y destino en las 8 tablas, rechazo efectivo de conexiones sin SSL
y acceso remoto directo de \texttt{postgres} bloqueado. Las consultas críticas
se ejecutan en el servidor PostgreSQL en 0.236 ms y 0.317 ms, muy por debajo de
cualquier umbral operativo para el volumen analizado. El proyecto se entrega
con scripts, evidencias y anexos.

\section{Descripción del cliente}

La Cooperativa de Ahorro y Crédito Horizonte PyME S.R.L. es una institución
financiera de alcance regional orientada a microcréditos para emprendedores,
comercios locales y pequeñas y medianas empresas. Su operación se organiza en
cinco sucursales (Central, Norte, Sur, Este y Oeste) y gira en torno a tres
procesos de negocio: la captación de ahorro, la colocación de crédito y el
seguimiento de la cartera. Los productos de crédito vigentes son personal,
empresarial, educativo, médico y agrícola, cada uno con montos máximos, tasas
y plazos diferenciados.

El sistema de información administra el ciclo de vida completo del socio:
registro y actualización de datos personales, apertura de cuentas de ahorro,
presentación y evaluación de solicitudes de crédito, desembolso, generación
de la tabla de amortización, registro de pagos y control de morosidad. Los
reportes principales son la cartera en mora (\texttt{vw\_cartera\_mora}) y el
resumen de saldos por socio (\texttt{vw\_resumen\_saldos\_socio}). La
información tratada es financiera y personal: nombres, documentos de identidad,
teléfonos, correos, direcciones, números de cuenta, saldos, montos de crédito
y historial de pagos.

Los interesados (\emph{stakeholders}) del proyecto son la Gerencia General,
que necesita continuidad operativa y información oportuna para la toma de
decisiones; la Unidad de Riesgos, que depende de la integridad de la cartera;
la Unidad de Auditoría, que requiere trazabilidad y acceso controlado; el
Área de Tecnología de la Información, responsable de la operación; y los
socios, cuyos datos deben permanecer confidenciales y disponibles. La
migración debe preservar la confianza de todos ellos: un dato perdido, una
cuenta mal conciliada o una filtración de identidad tendrían consecuencias
reputacionales, regulatorias y económicas.

\subsection{Requisitos del cliente}

A partir de las necesidades expresadas por los interesados se derivaron los
requisitos que el sistema destino debía satisfacer. La
Tabla~\ref{tab:requisitos-funcionales} enumera los requisitos funcionales.

\begin{table}[H]
  \centering
  \caption{Requisitos funcionales del cliente}
  \label{tab:requisitos-funcionales}
  \contenidoConFuente{%
    \begin{tablaAPA}
      \begin{tabular}{@{}C{0.9cm}L{4.4cm}L{10.2cm}@{}}
        \toprule
        \textbf{N.} & \textbf{Requisito} & \textbf{Descripción} \\
        \midrule
        1 & Migración íntegra & Traslado de las 8 tablas y las 2 vistas a PostgreSQL 16 preservando identificadores, integridad referencial y zona horaria. \\
        2 & Endurecimiento del destino & TLS obligatorio, autenticación \texttt{scram-sha-256}, roles de mínimo privilegio y registro de conexiones y cambios de estructura. \\
        3 & Enmascaramiento & Vista de auditoría con CI/NIT, teléfono, correo y cuenta enmascarados, sin exponer las tablas base. \\
        4 & Respaldo cifrado y restauración & Generación de un respaldo cifrado y prueba de su restauración completa. \\
        5 & Automatización y evidencias & Auditoría, migración y verificación con scripts reproducibles y evidencias reales. \\
        \bottomrule
      \end{tabular}
    \end{tablaAPA}%
  }{Elaboración propia.}
\end{table}

La Tabla~\ref{tab:requisitos-no-funcionales} complementa la lista anterior con
los requisitos no funcionales.

\begin{table}[H]
  \centering
  \caption{Requisitos no funcionales del cliente}
  \label{tab:requisitos-no-funcionales}
  \contenidoConFuente{%
    \begin{tablaAPA}
      \begin{tabular}{@{}C{0.9cm}L{4.4cm}L{10.2cm}@{}}
        \toprule
        \textbf{N.} & \textbf{Requisito} & \textbf{Descripción} \\
        \midrule
        1 & Infraestructura & Ejecución de todo el entorno con Docker Compose, sin instalar motores en el sistema anfitrión. \\
        2 & Sistema operativo y red & Linux en el laboratorio y puertos publicados solo en \texttt{127.0.0.1}. \\
        3 & Zona horaria & \texttt{America/La\_Paz} para fechas y reportes. \\
        4 & Entregables & Informe con la plantilla UPDS y demostración en vivo de 15--20 minutos, sin modificar \texttt{plantilla-latex/}. \\
        5 & Datos & Datos sintéticos, sin información personal real. \\
        \bottomrule
      \end{tabular}
    \end{tablaAPA}%
  }{Elaboración propia.}
\end{table}

Cada requisito se verifica en el Capítulo~IV y se resume en la lista de
cotejo del Anexo~C. La combinación de TLS, roles de mínimo privilegio,
enmascaramiento y respaldo cifrado exigida por el cliente se corresponde con
los controles recomendados por CIS, OWASP y NIST
\citep{cis2024postgres,owasp2023asvs,nist80053}.

\section{Problema de contexto}

MySQL 5.7 alcanzó el fin de su vida útil el 31 de octubre de 2023
\citep{oracle2023mysql57}. A partir de esa fecha Oracle dejó de publicar
parches de seguridad y correcciones para la rama 5.7, de modo que cualquier
vulnerabilidad descubierta queda sin remedio oficial. Para una entidad
financiera, operar sobre un motor sin soporte implica aceptar un riesgo
tecnológico y regulatorio que crece con el tiempo. El problema no es
únicamente la versión: el esquema heredado arrastra decisiones de diseño que
hoy resultan costosas.

En primer lugar, el usuario de aplicación \texttt{app\_legacy\_user} tiene
\texttt{GRANT ALL PRIVILEGES} sobre toda la base. Si la aplicación sufre una
inyección SQL o un compromiso de credenciales, el atacante hereda permisos de
lectura, escritura, eliminación y modificación de estructura. No existe
separación entre la cuenta que opera el sistema y la cuenta que administra la
base, ni entre la operación y la auditoría. En segundo lugar, las conexiones
no exigen cifrado: el tráfico entre la aplicación y el motor puede viajar en
texto claro dentro de la red, lo que expone credenciales y datos financieros.
En tercer lugar, el esquema no declara claves foráneas, de modo que la
integridad referencial depende por completo del código de la aplicación; una
carga manual, un \emph{script} defectuoso o un borrado accidental pueden dejar
cuentas sin socio, cuotas sin préstamo o transacciones sin cuenta.

A lo anterior se suman las incompatibilidades de tipos. \texttt{TINYINT(1)}
se usa como booleano, pero MySQL acepta valores distintos de 0 y 1;
\texttt{ENUM} restringe los estados en el motor, pero no se traduce de forma
directa a PostgreSQL; \texttt{DATETIME} no almacena zona horaria, lo que
introduce ambigüedad en operaciones distribuidas; \texttt{INT UNSIGNED}
duplica el rango positivo pero no existe como tipo nativo en PostgreSQL; y la
collation \texttt{utf8\_general\_ci} corresponde a \texttt{utf8mb3}, una
codificación de tres bytes que no cubre todos los caracteres Unicode. La
Tabla~\ref{tab:riesgos-origen} resume los riesgos identificados y su efecto
potencial.

\begin{table}[H]
  \centering
  \caption{Riesgos detectados en el sistema de origen}
  \label{tab:riesgos-origen}
  \contenidoConFuente{%
    \begin{tablaAPA}
      \begin{tabular}{@{}L{3.1cm}L{5.3cm}L{6.2cm}@{}}
        \toprule
        \textbf{Riesgo} & \textbf{Evidencia} & \textbf{Impacto potencial} \\
        \midrule
        Motor sin soporte & MySQL 5.7.44, EOL 31/10/2023 & Vulnerabilidades sin parche; incumplimiento de políticas de seguridad. \\
        Privilegios excesivos & \texttt{GRANT ALL} a \texttt{app\_legacy\_user} & Compromiso total de la base ante una inyección o robo de credenciales. \\
        Tráfico sin cifrado & Sin TLS obligatorio en el origen & Exposición de credenciales, saldos e identidades en la red. \\
        Ausencia de claves foráneas & 0 \texttt{FOREIGN KEY} en el esquema & Huérfanos, inconsistencias y reportes financieros erróneos. \\
        Tipos ambiguos & 7 \texttt{ENUM}, 10 \texttt{DATETIME}, 4 \texttt{TINYINT(1)} & Conversiones silenciosas, pérdida de zona horaria y validaciones débiles. \\
        Collation limitada & \texttt{utf8\_general\_ci} (\texttt{utf8mb3}) & Caracteres no representables y comparaciones inconsistentes. \\
        Auditoría insuficiente & Sin registro de conexiones ni DDL & Imposibilidad de reconstruir accesos y cambios. \\
        Respaldo no cifrado & \texttt{mysqldump} en texto claro & Exposición de datos si el archivo de respaldo se filtra. \\
        \bottomrule
      \end{tabular}
    \end{tablaAPA}%
  }{Elaboración propia con datos de \texttt{evidence/01\_inventario\_origen.txt}.}
\end{table}

La cooperativa necesita modernizar su plataforma de datos sin detener la
operación, sin perder información y sin repetir los defectos de seguridad del
sistema heredado. La migración debe ser verificable, reversible y reproducible
en cualquier equipo del grupo de trabajo.

\section{Solución propuesta}

La solución se apoya en cinco decisiones técnicas. La migración a
PostgreSQL 16 aporta un motor con soporte activo, cumplimiento ACID, tipos
nativos para booleanos, marcas de tiempo con zona horaria, identidades,
restricciones \texttt{CHECK}, claves foráneas e índices avanzados
\citep{postgresql2024}. El entorno completo se ejecuta con Docker Compose: un
servicio MySQL 5.7 para el origen, un servicio PostgreSQL 16 para el destino y
un contenedor migrador con Python, \texttt{pymysql} y \texttt{psycopg2}. Así se
evita instalar motores en el sistema anfitrión y se garantiza que cualquier
integrante levante el mismo escenario con un solo comando
\citep{docker2024,compose2024}.

Las transformaciones se definieron de forma explícita, regla por regla, en
lugar de depender de una conversión automática genérica. El volumen es pequeño
(0.28 MB), pero las reglas son específicas: convertir \texttt{TINYINT(1)} a
\texttt{BOOLEAN}, \texttt{ENUM} a \texttt{VARCHAR} con \texttt{CHECK},
\texttt{DATETIME} a \texttt{TIMESTAMPTZ} en la zona \texttt{America/La\_Paz},
preservar los identificadores originales con \texttt{OVERRIDING SYSTEM VALUE}
y reajustar las secuencias de identidad. Un ETL explícito permite auditar cada
regla y repetir la carga de forma determinista.

El proceso de migración se construyó sobre \textbf{pgloader}
con schema de staging (\texttt{stg}), validación SQL previa (\texttt{precheck}),
carga en una sola transacción con rollback garantizado (\texttt{postcheck}),
reajuste de secuencias y \texttt{ANALYZE}. El ETL artesanal en Python
(\texttt{src/migration/migrate.py}) fue archivado como legado y ya no se
ejecuta.

La seguridad se aplicó por capas. El perímetro exige TLS;
la autenticación usa \texttt{scram-sha-256}; la autorización se basa en roles
con privilegios mínimos; los datos sensibles se exponen a auditoría solo a
través de una vista enmascarada; las conexiones y el DDL se registran; y los
respaldos se cifran con AES-256-CBC. La Tabla~\ref{tab:solucion-componentes}
relaciona cada componente con su propósito y su evidencia.

\begin{table}[H]
  \centering
  \caption{Componentes de la solución propuesta}
  \label{tab:solucion-componentes}
  \contenidoConFuente{%
    \begin{tablaAPA}
      \begin{tabular}{@{}L{3.5cm}L{6.0cm}L{5.0cm}@{}}
        \toprule
        \textbf{Componente} & \textbf{Propósito} & \textbf{Evidencia} \\
        \midrule
        \texttt{docker-compose.yml} & Orquestar origen, destino, migrador y documentación & Servicios \emph{healthy}; puertos en \texttt{127.0.0.1}. \\
        \texttt{src/origin/} & Esquema y datos legacy reproducibles & 8 tablas, 2 vistas, 1\,469 filas. \\
        \texttt{src/target/} & Esquema destino, roles y configuración endurecida & 9 claves foráneas, 3 roles, SSL y \texttt{pg\_hba}. \\
        \texttt{src/migration/} & pgloader + staging, precheck, postcheck y rollback transaccional & 1\,469 filas, 0 rechazos, 0 errores. \\
        \texttt{scripts/} & Auditoría, migración, respaldo y pruebas & 3 scripts y 3 evidencias reales. \\
        \texttt{docs/} & Informe, anexos y guion de demostración & PDF compilable y guion de 15--20 min. \\
        \bottomrule
      \end{tabular}
    \end{tablaAPA}%
  }{Elaboración propia.}
\end{table}

\section{Alineación con los ODS 8 y 9}

El proyecto se alinea con el Objetivo de Desarrollo Sostenible 8 (Trabajo
decente y crecimiento económico) porque mejora la capacidad de la cooperativa
para otorgar microcréditos de forma oportuna. La meta 8.3 promueve la
formalización y el crecimiento de las micro, pequeñas y medianas empresas; la
meta 8.10 busca fortalecer la capacidad de las instituciones financieras para
ampliar el acceso a servicios bancarios y financieros \citep{un2023ods8}. Al
reducir el tiempo de aprobación y desembolso, y al disminuir el riesgo de
errores en la cartera, la migración contribuye directamente a que más
emprendedores accedan a capital de trabajo.

El ODS 9 (Industria, innovación e infraestructura) se relaciona con la
modernización de una infraestructura crítica de datos. La
meta 9.1 promueve infraestructuras fiables, sostenibles y resilientes; la meta
9.5 fortalece la investigación y la capacidad tecnológica; y la meta 9.c
busca ampliar el acceso a las tecnologías de la información y las
comunicaciones \citep{un2023ods9}. La contenerización, el cifrado, la
automatización de pruebas y la documentación reproducible son aportes
concretos a una infraestructura más resiliente y gobernable.

\begin{table}[H]
  \centering
  \caption{Contribución del proyecto a los ODS 8 y 9}
  \label{tab:ods}
  \contenidoConFuente{%
    \begin{tablaAPA}
      \begin{tabular}{@{}L{1.8cm}L{2.2cm}L{5.6cm}L{4.9cm}@{}}
        \toprule
        \textbf{ODS} & \textbf{Meta} & \textbf{Aporte del proyecto} & \textbf{Indicador de logro} \\
        \midrule
        ODS 8 & 8.3 & Agiliza la evaluación y el desembolso de microcréditos a pymes y emprendedores. & Tiempo de consulta de cartera y saldos en el orden de milisegundos. \\
        ODS 8 & 8.10 & Amplía el acceso a servicios financieros con información íntegra y disponible. & 55 socios y 66 cuentas migrados sin pérdida. \\
        ODS 9 & 9.1 & Sustituye una infraestructura obsoleta por una plataforma resiliente y reproducible. & Entorno Docker Compose levantado desde cero. \\
        ODS 9 & 9.5 & Incorpora automatización, pruebas y documentación técnica reutilizable. & 75 pruebas PASS y 3 evidencias reales. \\
        ODS 9 & 9.c & Facilita el acceso seguro a TIC mediante TLS, SCRAM y roles mínimos. & Conexiones sin SSL rechazadas; 7 roles auditados. \\
        \bottomrule
      \end{tabular}
    \end{tablaAPA}%
  }{Elaboración propia con base en \citet{un2023ods8} y \citet{un2023ods9}.}
\end{table}




% --- CAPÍTULO II -------------------------------------------------------------
\caratulaCapitulo{CAPÍTULO II}{OBJETIVOS, ALCANCE Y LÍMITES}
% ============================================================================
%  CAPÍTULO II: OBJETIVOS, ALCANCE Y LÍMITES
% ============================================================================

\section{Objetivos}

\subsection{Objetivo general}

Migrar de forma íntegra, segura y verificable la base de datos de créditos y
ahorros de la Cooperativa de Ahorro y Crédito Horizonte PyME S.R.L. desde
MySQL 5.7 hacia PostgreSQL 16, sobre un entorno Linux reproducible con Docker
Compose, aplicando endurecimiento, roles de mínimo privilegio, enmascaramiento
de datos sensibles y respaldos cifrados.

\subsection{Objetivos específicos}

\begin{enumerate}
  \item Diagnosticar el sistema de origen: versión, tamaño, tablas, vistas,
        índices, usuarios, privilegios y riesgos de seguridad.
  \item Identificar y documentar las incompatibilidades de tipos y semántica
        entre MySQL 5.7 y PostgreSQL 16, con su regla de conversión.
  \item Diseñar el esquema destino con tipos nativos, restricciones
        \texttt{CHECK}, claves foráneas, índices y vistas equivalentes.
  \item Implementar un ETL determinista que preserve identificadores,
        integridad y zona horaria, y que verifique conteos origen--destino.
  \item Endurecer PostgreSQL 16 con TLS/SSL obligatorio,
        \texttt{scram-sha-256}, revocación de \texttt{PUBLIC}, roles granulares
        y registro de conexiones y DDL.
  \item Enmascarar documentos, teléfonos, correos y números de cuenta en una
        vista de auditoría, sin exponer las tablas base.
  \item Generar respaldos cifrados con AES-256-CBC, calcular su huella
        SHA-256 y probar la restauración completa.
  \item Automatizar la auditoría, la migración y la verificación en scripts
        reproducibles, con evidencias reales en \texttt{evidence/}.
\end{enumerate}

\section{Alcance y límites}

\subsection{Alcance}

El alcance incluye el entorno de laboratorio completo: los servicios Docker
de origen y destino, el esquema y los datos sintéticos del origen, el esquema
destino, el ETL, los scripts de auditoría, migración, respaldo y verificación,
las evidencias generadas y la documentación técnica del proceso de migración
y endurecimiento. Se
cubren las 8 tablas, las 2 vistas, los 19 índices, el usuario legacy, las
incompatibilidades de tipos, el hardening del motor destino, los roles, el
enmascaramiento, el cifrado en tránsito y en reposo, y las pruebas de
integridad y rendimiento básico.

\subsection{Límites}

Quedan fuera del alcance la puesta en producción real, la migración de datos
históricos de la cooperativa, la modificación del código de la aplicación, la
implementación de alta disponibilidad con réplicas o \emph{failover}, la
instalación de un \emph{pool} de conexiones como PgBouncer, la certificación
formal de cumplimiento normativo y la operación 24×7. Los datos utilizados
son ficticios y deterministas; no se manipuló información real de personas.
El entorno se ejecuta en un solo host Linux y no representa una arquitectura
de producción multizona. Estas exclusiones se documentan para que el lector
no confunda el laboratorio con un despliegue productivo.


% --- CAPÍTULO III ------------------------------------------------------------
\caratulaCapitulo{CAPÍTULO III}{DIAGNÓSTICO DEL ORIGEN Y HARDENING DEL DESTINO}
% ============================================================================
%  CAPÍTULO II: DIAGNÓSTICO DEL ORIGEN Y HARDENING DEL DESTINO
% ============================================================================

\section{Arquitectura del sistema de origen}

El sistema de origen es una instalación MySQL 5.7.44 sobre Linux, ejecutada
en el contenedor \texttt{coop\_horizonte\_legacy}. La base se denomina
\texttt{coop\_horizonte\_legacy}, usa el motor InnoDB y la collation
\texttt{utf8\_general\_ci}. La aplicación legacy se conecta con el usuario
\texttt{app\_legacy\_user}, creado con \texttt{GRANT ALL PRIVILEGES} sobre
toda la base. No existe una cuenta separada para auditoría, no se exige TLS
y no hay un esquema de respaldo cifrado dentro del flujo operativo.

La Figura~\ref{fig:arquitectura-origen} representa la arquitectura
encontrada. La aplicación y el motor comparten la red del contenedor; el
tráfico no está cifrado y el usuario de aplicación puede ejecutar DDL, DML y
DCL. El puerto 3306 se publica únicamente en \texttt{127.0.0.1}, lo que
reduce la exposición desde la red del anfitrión, pero no corrige los
problemas internos de privilegios ni la falta de cifrado.

\figuraAPA{fig:arquitectura-origen}{Arquitectura del sistema de origen MySQL 5.7}{%
  \begin{tikzpicture}[
      nodo/.style={draw, rounded corners, minimum width=3.2cm,
                   minimum height=1.0cm, align=center, font=\small},
      riesgo/.style={draw, dashed, rounded corners, minimum width=3.2cm,
                     minimum height=0.8cm, align=center, font=\scriptsize,
                     text=black!70},
      linea/.style={-Stealth, thick}]
    \node[nodo] (app) {Aplicación\\legacy};
    \node[nodo, right=2.4cm of app] (db) {MySQL 5.7.44\\\texttt{coop\_horizonte\_legacy}};
    \node[riesgo, below=0.9cm of app] (r1) {Sin TLS};
    \node[riesgo, below=0.9cm of db] (r2) {\texttt{GRANT ALL}};
    \draw[linea] (app) -- node[above, font=\scriptsize]{TCP/3306} (db);
    \draw[dashed, red!60] (r1) -- (app);
    \draw[dashed, red!60] (r2) -- (db);
  \end{tikzpicture}%
}{Elaboración propia con datos de \texttt{evidence/01\_inventario\_origen.txt}.}

\capturaAPA{fig:cap-origen-version}{Versión y base de datos del sistema de origen}%
  {captura-origen-version.png}%
  {Salida de \texttt{SELECT VERSION();} y \texttt{SHOW DATABASES;} en MySQL 5.7.44.}

El origen contiene 8 tablas de negocio, 2 vistas, 19 entradas de índice y un
solo usuario de aplicación. La Tabla~\ref{tab:inventario-origen} resume el
inventario. El tamaño total reportado por \texttt{information\_schema} es de
0.28 MB, un volumen pequeño que permite una ventana de migración corta y un
enfoque de corte controlado (\emph{big bang}) en lugar de replicación continua.

\begin{table}[H]
  \centering
  \caption{Inventario de objetos del sistema de origen}
  \label{tab:inventario-origen}
  \contenidoConFuente{%
    \begin{tablaAPA}
      \begin{tabular}{@{}L{4.6cm}C{2.0cm}L{7.9cm}@{}}
        \toprule
        \textbf{Tipo de objeto} & \textbf{Cantidad} & \textbf{Detalle} \\
        \midrule
        Tablas base & 8 & \texttt{sucursales}, \texttt{socios}, \texttt{cuentas\_ahorro}, \texttt{tipos\_credito}, \texttt{solicitudes\_credito}, \texttt{prestamos}, \texttt{cuotas\_amortizacion}, \texttt{transacciones}. \\
        Vistas & 2 & \texttt{vw\_cartera\_mora} y \texttt{vw\_resumen\_saldos\_socio}. \\
        Entradas de índice & 19 & 8 claves primarias, 3 \texttt{UNIQUE} y 8 índices no únicos (uno compuesto). \\
        Secuencias & 0 & MySQL 5.7 no define objetos \texttt{SEQUENCE}; 8 columnas \texttt{AUTO\_INCREMENT} pasan a identidades en PostgreSQL. \\
        Usuarios de aplicación & 1 & \texttt{app\_legacy\_user} con \texttt{GRANT ALL}. \\
        Claves foráneas & 0 & La integridad referencial no está declarada en el motor. \\
        Tamaño de la base & 0.28 MB & Medido con \texttt{data\_length + index\_length}. \\
        \bottomrule
      \end{tabular}
    \end{tablaAPA}%
  }{Elaboración propia con datos de \texttt{evidence/01\_inventario\_origen.txt}.}
\end{table}

La ausencia de objetos \texttt{SEQUENCE} no implica que el destino carezca de
generadores de identificadores: las 8 columnas \texttt{AUTO\_INCREMENT} se
convierten en identidades \texttt{BIGINT GENERATED ALWAYS AS IDENTITY} y sus
valores se reajustan con \texttt{setval} al final de la carga, sin inventar
objetos nuevos ni perder los identificadores originales. {\hyphenation{OVERRIDING}%
La documentación oficial de PostgreSQL 16 describe la identidad así: «In an
INSERT command, if ALWAYS is selected, a user-specified value is only accepted
if the INSERT statement specifies OVERRIDING SYSTEM VALUE»
\citep[apartado «GENERATED ALWAYS AS IDENTITY»]{postgresql2024}.}

\capturaAPA{fig:cap-inventario}{Inventario de objetos del sistema de origen}%
  {captura-inventario-objetos.png}%
  {Salida de \texttt{scripts/01\_audit\_origin.sh} con tablas, vistas, índices, usuarios y privilegios.}

La salida completa del inventario se presenta en el Anexo~C
(Figuras~\ref{fig:anexo-inventario-1} y~\ref{fig:anexo-inventario-2}).

\subsection{Diagrama de clases de la base de datos origen}

La Figura~\ref{fig:diagrama-clases-origen} presenta el diagrama de clases de
la base de datos origen. Cada clase representa una tabla de
\texttt{coop\_horizonte\_legacy}, con sus atributos y tipos de datos de
MySQL 5.7. Las asociaciones reflejan las relaciones lógicas del negocio: un
socio tiene muchas cuentas de ahorro y muchas solicitudes de crédito; una
solicitud aprobada genera un préstamo; un préstamo tiene muchas cuotas de
amortización; y las transacciones referencian un socio, una cuenta o un
préstamo según el tipo de movimiento. Estas relaciones son lógicas, inferidas
del modelo de datos y del código de la aplicación: como se documenta en la
Sección~\ref{sec:riesgos}, el origen no declara claves foráneas en el motor,
y solo se formalizan como restricciones \texttt{FOREIGN KEY} en el destino.

\noindent\begin{minipage}{\linewidth}
\capturaAPA{fig:diagrama-clases-origen}{Diagrama de clases de la base de datos origen}%
  {diagrama-clases-origen.png}%
  {Elaboración propia con base en \texttt{src/origin/01\_schema\_origin.sql}.}
\end{minipage}

\section{Inventario de datos y volumen}

El conjunto de datos fue generado de forma determinista con una semilla fija
(\texttt{seed=42}) para que cualquier integrante obtenga exactamente los
mismos valores. La Tabla~\ref{tab:conteos-origen} presenta los conteos
exactos. Destacan los 55 socios, las 66 cuentas de ahorro (algunos socios
tienen dos cuentas), las 50 solicitudes de crédito, los 32 préstamos
aprobados, las 516 cuotas de amortización y las 740 transacciones. El total
es de 1\,469 filas.

\begin{table}[H]
  \centering
  \caption{Conteo de filas en el sistema de origen}
  \label{tab:conteos-origen}
  \contenidoConFuente{%
    \begin{tablaAPA}
      \begin{tabular}{@{}L{6.2cm}R{3.0cm}L{5.0cm}@{}}
        \toprule
        \textbf{Tabla} & \textbf{Filas} & \textbf{Propósito} \\
        \midrule
        \texttt{sucursales} & 5 & Catálogo de agencias. \\
        \texttt{socios} & 55 & Datos personales y estado del socio. \\
        \texttt{cuentas\_ahorro} & 66 & Cuentas y saldos iniciales. \\
        \texttt{tipos\_credito} & 5 & Productos de crédito. \\
        \texttt{solicitudes\_credito} & 50 & Solicitudes y su estado. \\
        \texttt{prestamos} & 32 & Créditos desembolsados. \\
        \texttt{cuotas\_amortizacion} & 516 & Plan de pagos. \\
        \texttt{transacciones} & 740 & Depósitos, retiros, pagos y desembolsos. \\
        \midrule
        \textbf{Total} & \textbf{1\,469} & 8 tablas de negocio. \\
        \bottomrule
      \end{tabular}
    \end{tablaAPA}%
  }{Elaboración propia con datos de \texttt{evidence/01\_inventario\_origen.txt} y \texttt{evidence/02\_ejecucion\_migracion.txt}.}
\end{table}

La distribución de filas no es arbitraria: las transacciones se derivan de
los desembolsos, los pagos de cuotas y los movimientos de las cuentas; las
cuotas se generan a partir de cada préstamo; y los préstamos solo existen
para solicitudes aprobadas. Esa coherencia es la que permite verificar la
integridad referencial después de la migración. El \emph{script}
\texttt{src/origin/generate\_seed.py} documenta las reglas de generación y
puede regenerar el archivo \texttt{02\_seed\_origin.sql} en cualquier momento.

\section{Riesgos de seguridad detectados}
\label{sec:riesgos}

\subsection{Privilegios excesivos del usuario de aplicación}

El hallazgo más grave es que \texttt{app\_legacy\_user} tiene
\texttt{GRANT ALL PRIVILEGES} sobre \texttt{coop\_horizonte\_legacy.*}. En
MySQL, \texttt{ALL} incluye \texttt{SELECT}, \texttt{INSERT}, \texttt{UPDATE},
\texttt{DELETE}, \texttt{CREATE}, \texttt{DROP}, \texttt{ALTER},
\texttt{INDEX}, \texttt{REFERENCES}, \texttt{CREATE TEMPORARY TABLES},
\texttt{LOCK TABLES}, \texttt{EXECUTE}, \texttt{CREATE VIEW},
\texttt{SHOW VIEW}, \texttt{CREATE ROUTINE}, \texttt{ALTER ROUTINE},
\texttt{EVENT} y \texttt{TRIGGER}. La aplicación no necesita DDL ni
\texttt{DROP}; sin embargo, si es comprometida, el atacante puede borrar
tablas, alterar el esquema o crear disparadores maliciosos. La mitigación es
separar la operación de la administración y otorgar solo DML sobre las tablas
necesarias, que es exactamente lo que hace el rol \texttt{app\_core\_rw} en
el destino.

\capturaAPA{fig:cap-privilegios}{Privilegios del usuario de aplicación en el origen}%
  {captura-privilegios-app.png}%
  {Salida de \texttt{SHOW GRANTS FOR 'app\_legacy\_user'@'\%';} con \texttt{GRANT ALL PRIVILEGES}.}

\subsection{Transporte sin cifrado}

MySQL 5.7 puede habilitar TLS, pero la instalación analizada no lo exige. La
aplicación se conecta sin \texttt{REQUIRE SSL}, de modo que las credenciales
y los datos pueden ser interceptados en la red interna. En el destino se
exige \texttt{hostssl} para todas las conexiones TCP y se rechaza cualquier
intento sin cifrado, como se demuestra en la
Sección~\ref{sec:evidencias-hardening}.

\subsection{Ausencia de integridad referencial}

El esquema no declara claves foráneas. La aplicación es la única responsable
de que \texttt{cuentas\_ahorro.socio\_id} apunte a un socio existente, que
\texttt{prestamos.solicitud\_id} corresponda a una solicitud aprobada y que
\texttt{cuotas\_amortizacion.prestamo\_id} pertenezca a un préstamo. Cualquier
error de código, carga manual o borrado accidental puede producir huérfanos.
En el destino se añadieron 9 claves foráneas y se verificó que no existieran
huérfanos después de la carga.

\subsection{Auditoría insuficiente}

No hay registro sistemático de conexiones, desconexiones ni cambios de
estructura. Si ocurriera un acceso no autorizado o un cambio de esquema, la
cooperativa no tendría evidencia para reconstruir los hechos. El destino
habilita \texttt{log\_connections}, \texttt{log\_disconnections} y
\texttt{log\_statement = 'ddl'}, y además separa los roles para que la cuenta
de auditoría no pueda modificar datos.

\subsection{Respaldo sin cifrado y sin prueba de restauración}

Un respaldo en texto claro es tan sensible como la base original. Si el
archivo se copia a un medio externo, se comparte por error o se almacena en
un servidor comprometido, los datos personales y financieros quedan
expuestos. El proyecto genera el respaldo del destino con
\texttt{pg\_dump -Fc} y lo cifra con AES-256-CBC y PBKDF2; además, descifra y
restaura el respaldo en una base temporal para comprobar que es utilizable.
Un respaldo que nunca se restaura no es un respaldo: es una esperanza.

\subsection{Fin de vida del motor}

MySQL 5.7 dejó de recibir parches oficiales en octubre de 2023. La
cooperativa no puede depender de un motor sin soporte para custodiar datos
financieros. PostgreSQL 16 tiene soporte activo, una comunidad amplia y un
ciclo de versiones predecible \citep{postgresql2024}. La migración reduce el
riesgo de vulnerabilidades no corregidas y facilita la aplicación de futuras
actualizaciones.

\section{Matriz de incompatibilidades y reglas de conversión}

La migración heterogénea no consiste en copiar valores: exige interpretar la
semántica de cada tipo. La Tabla~\ref{tab:incompatibilidades} presenta las
incompatibilidades reales encontradas, su frecuencia y la regla aplicada en el
destino. Las decisiones se justifican con la documentación oficial de
PostgreSQL \citep{postgresql2024} y con los principios de diseño relacional
\citep{elmasri2016,silberschatz2020}.

\begin{table}[H]
  \centering
  \caption{Matriz de incompatibilidades MySQL 5.7--PostgreSQL 16}
  \label{tab:incompatibilidades}
  \contenidoConFuente{%
    \begin{tablaAPA}
      \begin{tabular}{@{}L{2.7cm}L{3.3cm}L{3.6cm}L{5.2cm}@{}}
        \toprule
        \textbf{MySQL 5.7} & \textbf{Columnas reales} & \textbf{PostgreSQL 16} & \textbf{Regla de conversión} \\
        \midrule
        \texttt{TINYINT(1)} & 4 columnas: \texttt{estatus} en \texttt{sucursales}, \texttt{socios} y \texttt{cuentas\_ahorro}; \texttt{pagado\_total} en \texttt{prestamos}. & \texttt{BOOLEAN} & Convertir 0 a \texttt{FALSE} y cualquier valor distinto de 0 a \texttt{TRUE}. Validar que no existan valores intermedios. \\
        \texttt{ENUM} & 7 columnas: estados de socios, tipos de crédito, solicitudes, préstamos, cuotas y transacciones. & \texttt{VARCHAR(n)} + \texttt{CHECK} & Reemplazar el tipo por texto con una restricción \texttt{CHECK ... IN (...)} que conserve el dominio. \\
        \texttt{DATETIME} & 10 columnas: fechas de creación, apertura, cierre, solicitud, aprobación, desembolso, pago y transacción. & \texttt{TIMESTAMPTZ} & Interpretar el valor original en \texttt{America/La\_Paz} y almacenarlo como instante con zona horaria. \\
        \texttt{INT UNSIGNED AUTO\_INCREMENT} & 8 claves primarias. & \texttt{BIGINT GENERATED ALWAYS AS IDENTITY} & Crear identidad, insertar los IDs originales con \texttt{OVERRIDING SYSTEM VALUE} y reajustar la secuencia con \texttt{setval}. \\
        \texttt{utf8\_general\_ci} & Base y tablas. & UTF-8 nativo & Usar la codificación UTF-8 del clúster; no se requiere collation por columna. \\
        \texttt{UNIQUE KEY} & 3 columnas únicas: \texttt{numero\_identificacion}, \texttt{num\_cuenta}, \texttt{solicitud\_id}. & \texttt{UNIQUE} & Conservar la restricción y crear el índice único asociado. \\
        Sin \texttt{FOREIGN KEY} & 0 restricciones en el origen. & 9 \texttt{FOREIGN KEY} & Declarar las relaciones y validar que no existan huérfanos antes de la carga. \\
        \texttt{DECIMAL(12,2)} & Montos, saldos, tasas y cuotas. & \texttt{NUMERIC(12,2)} & Mantener precisión y escala; añadir \texttt{CHECK} para montos no negativos. \\
        \texttt{TEXT} & Descripciones, observaciones y direcciones. & \texttt{TEXT} & Conservar sin cambios; PostgreSQL no requiere prefijos de longitud. \\
        \bottomrule
      \end{tabular}
    \end{tablaAPA}%
  }{Elaboración propia con datos de \texttt{evidence/01\_inventario\_origen.txt} y \texttt{src/target/01\_schema\_target.sql}.}
\end{table}

\subsection{Incompatibilidades de sintaxis y adaptación}

Además de los tipos, la sintaxis de ambos motores difiere. La adaptación no
se resolvió con una conversión automática: se revisaron las sentencias del
origen y se aplicaron las equivalencias de la
Tabla~\ref{tab:incompatibilidades-sintaxis}.

\begin{table}[H]
  \centering
  \caption{Incompatibilidades de sintaxis MySQL 5.7--PostgreSQL 16}
  \label{tab:incompatibilidades-sintaxis}
  \contenidoConFuente{%
    \begin{tablaAPA}
      \begin{tabular}{@{}L{4.0cm}L{4.3cm}L{6.5cm}@{}}
        \toprule
        \textbf{MySQL 5.7} & \textbf{PostgreSQL 16} & \textbf{Decisión} \\
        \midrule
        \code{ENGINE=InnoDB}, \code{CHARSET}, \code{COLLATE} & Sintaxis nativa & Se omiten; se usa la codificación UTF-8 del clúster. \\
        \code{AUTO\_INCREMENT} & \code{GENERATED ALWAYS AS IDENTITY} & Se conservan los IDs y se reajusta con \texttt{setval}. \\
        \code{GRANT ALL ON db.*} con comillas invertidas & Roles y \texttt{GRANT} sobre el esquema & Se separan \texttt{admin\_dba}, \texttt{app\_core\_rw} y \texttt{auditor\_ro}. \\
        \code{SHOW GRANTS} & \texttt{pg\_roles} y \texttt{has\_table\_privilege} & Los privilegios se auditan por catálogo. \\
        \code{IFNULL} & \code{COALESCE} & Se adopta la función estándar SQL. \\
        \code{GROUP\_CONCAT} & \code{string\_agg} & Se usa en las vistas de reporte. \\
        \code{ON DUPLICATE KEY UPDATE} & \code{INSERT ... ON CONFLICT} & Disponible para cargas tipo \emph{upsert}. \\
        Comentarios con \texttt{\#} & Comentarios con \texttt{-{}-} & Se convierten al adaptar los scripts. \\
        Comillas invertidas para identificadores & Comillas dobles o sin delimitar & Se prefiere SQL estándar. \\
        \bottomrule
      \end{tabular}
    \end{tablaAPA}%
  }{Elaboración propia con base en \texttt{src/origin/01\_schema\_origin.sql} y \texttt{src/target/01\_schema\_target.sql}.}
\end{table}

Estas diferencias explican por qué el DDL y el ETL se escribieron
explícitamente en lugar de generarse con una herramienta genérica.

\subsection{Booleanos: de \texttt{TINYINT(1)} a \texttt{BOOLEAN}}

MySQL no tiene un tipo booleano nativo: \texttt{BOOLEAN} es un alias de
\texttt{TINYINT(1)} y acepta cualquier valor de 8 bits. PostgreSQL tiene
\texttt{BOOLEAN} con valores \texttt{TRUE}, \texttt{FALSE} y \texttt{NULL}.
La conversión aplicada en el ETL es explícita: \texttt{bool(int(valor))}. Si
el origen contuviera un 2 o un -1, el ETL lo interpretaría como verdadero; por
eso se recomienda una validación previa en producción. En el conjunto de
datos analizado todos los valores son 0 o 1.

\subsection{Estados: de \texttt{ENUM} a \texttt{VARCHAR} con \texttt{CHECK}}

\texttt{ENUM} es una enumeración almacenada como entero interno. Su principal
ventaja es la validación en el motor; su principal desventaja es que añadir un
valor requiere \texttt{ALTER TABLE} y que el orden interno puede cambiar si se
reordena la lista. En PostgreSQL se optó por \texttt{VARCHAR(10)} con una
restricción \texttt{CHECK}. Así se conserva la validación, se facilita la
lectura y se evita depender de un tipo propietario. Las restricciones se
nombraron de forma descriptiva para que los errores sean comprensibles.

\subsection{Fechas: de \texttt{DATETIME} a \texttt{TIMESTAMPTZ}}

\texttt{DATETIME} almacena una fecha y hora sin zona. Eso significa que el
valor \texttt{2026-11-20 08:09:40} no indica si ocurrió en La Paz, en UTC o en
otra zona. Para una entidad financiera, la ambigüedad es inaceptable: un pago
registrado a las 23:30 puede pertenecer a un día contable distinto según la
zona. La migración interpreta cada valor en \texttt{America/La\_Paz} y lo
almacena como \texttt{TIMESTAMPTZ}. El ETL usa \texttt{ZoneInfo} de Python
para asignar la zona antes de insertar, y el servidor configura
\texttt{timezone = 'America/La\_Paz'} para que las consultas y los reportes
se muestren en la hora local.

\subsection{Identificadores: de \texttt{AUTO\_INCREMENT} a identidad}

MySQL usa \texttt{AUTO\_INCREMENT} con enteros sin signo. PostgreSQL 16 usa
identidades (\texttt{GENERATED ALWAYS AS IDENTITY}), que son estándar SQL y
permiten controlar el ciclo de vida de la secuencia. El ETL inserta los IDs
originales con \texttt{OVERRIDING SYSTEM VALUE} para no romper las referencias
y, al final, ejecuta \texttt{setval} sobre cada secuencia. Sin ese ajuste, el
primer \texttt{INSERT} nuevo intentaría usar un ID ya ocupado y fallaría.

\subsection{Collation y codificación}

\texttt{utf8\_general\_ci} es la collation de \texttt{utf8mb3}, una
codificación de tres bytes que no puede representar todos los caracteres
Unicode (por ejemplo, emojis o ciertos símbolos). PostgreSQL usa UTF-8 nativo
y una collation del clúster. La migración no requiere conversión de datos
porque los valores del origen son ASCII y español; sin embargo, el cambio
elimina la limitación y prepara la base para nombres y direcciones con
caracteres adicionales.

\subsection{Integridad referencial}

La ausencia de claves foráneas en el origen obligó a validar la integridad
antes de declararlas en el destino. El ETL carga en orden jerárquico
(\texttt{sucursales}, \texttt{tipos\_credito}, \texttt{socios},
\texttt{cuentas\_ahorro}, \texttt{solicitudes\_credito}, \texttt{prestamos},
\texttt{cuotas\_amortizacion}, \texttt{transacciones}) y luego las pruebas
consultan cada relación con \texttt{LEFT JOIN} para contar huérfanos. El
resultado fue cero en las 9 relaciones, como se muestra en el
Capítulo~IV.

% ============================================================================
%  CAPÍTULO II (continuación): HARDENING DEL DESTINO
% ============================================================================

\section{Hardening del SGBD de destino}

El endurecimiento de PostgreSQL 16 se diseñó con un principio rector:
\emph{negar por defecto, conceder por excepción}. Cada capa ---transporte,
autenticación, autorización, datos y auditoría--- tiene una configuración
explícita y una prueba que la verifica. No se asumió que una medida estaba
aplicada porque el archivo de configuración la mencionara: se consultó el
estado efectivo del servidor y se registró la salida en
\texttt{evidence/03\_pruebas\_y\_hardening.txt}. La Tabla~\ref{tab:capas-hardening}
resume las capas y su propósito.

\begin{table}[H]
  \centering
  \caption{Capas de seguridad aplicadas en PostgreSQL 16}
  \label{tab:capas-hardening}
  \contenidoConFuente{%
    \begin{tablaAPA}
      \begin{tabular}{@{}L{2.7cm}L{5.2cm}L{6.4cm}@{}}
        \toprule
        \textbf{Capa} & \textbf{Medida} & \textbf{Propósito} \\
        \midrule
        Transporte & TLS 1.2 o superior, cifrados \texttt{HIGH}, \texttt{hostssl} obligatorio. & Impedir que credenciales y datos viajen en texto claro. \\
        Autenticación & \texttt{scram-sha-256} para todos los roles de aplicación. & Evitar el envío de contraseñas en claro y el uso de hashes débiles. \\
        Autorización & Revocación de \texttt{PUBLIC}; roles \texttt{admin\_dba}, \texttt{app\_core\_rw} y \texttt{auditor\_ro}. & Aplicar mínimo privilegio y separación de funciones. \\
        Datos & Vista \texttt{vw\_socios\_audit} con enmascaramiento de CI/NIT, teléfono, correo y cuenta. & Permitir auditoría sin exponer identificadores completos. \\
        Auditoría & \texttt{log\_connections}, \texttt{log\_disconnections}, \texttt{log\_statement = 'ddl'}. & Reconstruir accesos y cambios de estructura. \\
        Respaldo & \texttt{pg\_dump -Fc} + AES-256-CBC + PBKDF2. & Proteger la confidencialidad de las copias. \\
        \bottomrule
      \end{tabular}
    \end{tablaAPA}%
  }{Elaboración propia con base en \citet{cis2024postgres}, \citet{owasp2023asvs} y \citet{nist80053}.}
\end{table}

\subsection{TLS/SSL obligatorio}

El servidor destino se configuró con \texttt{ssl = on}, protocolo mínimo
TLS 1.2 y cifrados \texttt{HIGH:!aNULL:!MD5}. La clave privada
\texttt{server.key} tiene permisos \texttt{600} y pertenece al usuario
\texttt{postgres}; el certificado \texttt{server.crt} tiene permisos
\texttt{644}. El certificado es autofirmado y contiene el nombre
\texttt{db-destino} en el campo \texttt{Subject Alternative Name}, además de
\texttt{localhost} y \texttt{127.0.0.1}. Es válido desde el 20 de noviembre de
2026 hasta el 22 de febrero de 2029. La Tabla~\ref{tab:certificado} resume sus
atributos.

\begin{table}[H]
  \centering
  \caption{Atributos del certificado TLS del destino}
  \label{tab:certificado}
  \contenidoConFuente{%
    \begin{tablaAPA}
      \begin{tabular}{@{}L{4.2cm}L{9.8cm}@{}}
        \toprule
        \textbf{Atributo} & \textbf{Valor} \\
        \midrule
        Sujeto & \texttt{C=BO, ST=La Paz, L=La Paz, O=Cooperativa Horizonte PyME S.R.L., OU=TI, CN=db-destino}. \\
        SAN & \texttt{DNS:db-destino}, \texttt{DNS:localhost}, \texttt{IP:127.0.0.1}. \\
        Validez & 20/11/2026 -- 22/02/2029. \\
        Uso de clave & \texttt{digitalSignature}, \texttt{keyEncipherment}, \texttt{keyCertSign}. \\
        Propósito extendido & \texttt{serverAuth}. \\
        Restricción básica & \texttt{CA:TRUE} (certificado autofirmado de laboratorio). \\
        Permisos & \texttt{server.key}: \texttt{600}; \texttt{server.crt}: \texttt{644}. \\
        \bottomrule
      \end{tabular}
    \end{tablaAPA}%
  }{Elaboración propia con datos de \texttt{openssl x509} y \texttt{src/target/ssl/}.}
\end{table}

El fragmento relevante de \texttt{src/target/postgresql.conf} es el siguiente.
La configuración se monta en el contenedor y se aplica al iniciar; no se
edita el archivo dentro del volumen de datos, de modo que la configuración
queda versionada junto con el código.

\begin{lstlisting}[language=bash,caption={Configuración TLS y de auditoría en \texttt{postgresql.conf}},label={lst:postgresql-conf}]
ssl = on
ssl_cert_file = '/var/lib/postgresql/ssl/server.crt'
ssl_key_file = '/var/lib/postgresql/ssl/server.key'
ssl_min_protocol_version = 'TLSv1.2'
ssl_ciphers = 'HIGH:!aNULL:!MD5'
ssl_prefer_server_ciphers = on

password_encryption = 'scram-sha-256'

log_connections = on
log_disconnections = on
log_statement = 'ddl'
log_line_prefix = '%m [%p] %q%u@%d '

timezone = 'America/La_Paz'
log_timezone = 'America/La_Paz'
\end{lstlisting}

La prueba de conexión se realizó desde el contenedor migrador con
\texttt{sslmode=require} y \texttt{sslmode=verify-ca}. En el segundo caso se
pasó \texttt{sslrootcert=/etc/postgresql/ssl/server.crt} para validar la
cadena. Ambas conexiones funcionaron. La conexión con \texttt{sslmode=disable}
fue rechazada con el mensaje \texttt{no pg\_hba.conf entry ... no encryption},
lo que demuestra que la política no es una recomendación: es una condición
de acceso.

\capturaAPA{fig:cap-ssl}{Verificación TLS: conexión rechazada sin SSL y aceptada con \texttt{verify-ca}}%
  {captura-ssl-verificacion.png}%
  {Conexión con \texttt{sslmode=disable} rechazada por \texttt{pg\_hba.conf} y conexión con \texttt{sslmode=verify-ca} aceptada.}

\subsection{Autenticación \texttt{scram-sha-256} y reglas \texttt{pg\_hba.conf}}

PostgreSQL 16 usa \texttt{scram-sha-256} como método de autenticación por
defecto para contraseñas nuevas. SCRAM es un mecanismo de desafío-respuesta
que no transmite la contraseña en claro y resiste ataques de repetición
\citep{rfc5802,rfc7677}. El archivo \texttt{pg\_hba.conf} se diseñó para que
solo las conexiones TCP cifradas puedan autenticarse, y para que
\texttt{postgres} no tenga acceso remoto. La Tabla~\ref{tab:pg-hba} reproduce
las reglas efectivas.

\begin{table}[H]
  \centering
  \caption{Reglas efectivas de \texttt{pg\_hba.conf}}
  \label{tab:pg-hba}
  \contenidoConFuente{%
    \begin{tablaAPA}
      \begin{tabular}{@{}L{1.5cm}L{3.4cm}L{3.0cm}L{3.2cm}L{3.3cm}@{}}
        \toprule
        \textbf{Tipo} & \textbf{Base} & \textbf{Usuario} & \textbf{Dirección} & \textbf{Método} \\
        \midrule
        \texttt{local} & todas & todos & socket local & \texttt{peer} \\
        \texttt{local} & todas & \texttt{postgres} & socket local & \texttt{peer} \\
        \texttt{hostssl} & \texttt{coop\_horizonte\_core} & \texttt{admin\_dba} & \texttt{0.0.0.0/0}, \texttt{::/0} & \texttt{scram-sha-256} \\
        \texttt{hostssl} & \texttt{coop\_horizonte\_core} & \texttt{app\_core\_rw} & \texttt{0.0.0.0/0}, \texttt{::/0} & \texttt{scram-sha-256} \\
        \texttt{hostssl} & \texttt{coop\_horizonte\_core} & \texttt{auditor\_ro} & \texttt{0.0.0.0/0}, \texttt{::/0} & \texttt{scram-sha-256} \\
        \texttt{hostssl} & \texttt{coop\_horizonte\_core} & \texttt{migration\_loader} & \texttt{0.0.0.0/0}, \texttt{::/0} & \texttt{scram-sha-256} \\
        \texttt{hostssl} & \texttt{coop\_horizonte\_core} & \texttt{nicolas}, \texttt{joaquin}, \texttt{david}, \texttt{gustavo} & \texttt{0.0.0.0/0}, \texttt{::/0} & \texttt{scram-sha-256} \\
        \bottomrule
      \end{tabular}
    \end{tablaAPA}%
  }{Elaboración propia con datos de \texttt{evidence/03\_pruebas\_y\_hardening.txt}.}
\end{table}

La ausencia de una regla \texttt{host} para \texttt{postgres} es deliberada.
El superusuario solo puede conectarse por socket local con autenticación
\texttt{peer}, lo que exige que el usuario del sistema operativo coincida con
el usuario de la base. En el contenedor, los scripts administrativos se
ejecutan con \texttt{docker compose exec -u postgres}, de modo que no se
necesita exponer \texttt{postgres} por TCP. Si un atacante obtuviera la
contraseña de \texttt{postgres}, no podría usarla remotamente.

\subsection{Revocación de \texttt{PUBLIC} y roles de mínimo privilegio}

Por defecto, PostgreSQL permite que \texttt{PUBLIC} tenga \texttt{USAGE} y
\texttt{CREATE} sobre el esquema \texttt{public}, y \texttt{CONNECT} sobre
cualquier base. Eso es incompatible con un entorno financiero. El script
\texttt{src/target/02\_hardening\_roles.sql} revoca esos privilegios y crea
siete roles con responsabilidades separadas: tres roles de operación
(\texttt{admin\_dba}, \texttt{app\_core\_rw}, \texttt{auditor\_ro}), un rol de
carga de datos (\texttt{migration\_loader}) y cuatro roles de administradores
del equipo (\texttt{nicolas}, \texttt{joaquin}, \texttt{david},
\texttt{gustavo}). La Tabla~\ref{tab:roles} presenta la matriz de privilegios.

\begin{table}[H]
  \centering
  \caption{Roles y privilegios en PostgreSQL 16}
  \label{tab:roles}
  \contenidoConFuente{%
    \begin{tablaAPA}
      \begin{tabular}{@{}L{2.5cm}L{2.3cm}L{4.0cm}L{5.6cm}@{}}
        \toprule
        \textbf{Rol} & \textbf{Atributos} & \textbf{Permite} & \textbf{Prohíbe} \\
        \midrule
        \texttt{admin\_dba} & \texttt{LOGIN}, \texttt{NOSUPERUSER}, \texttt{CREATEDB}, \texttt{CREATEROLE}, \texttt{NOINHERIT} & Administrar la base, crear objetos, ejecutar DDL y DML, gestionar roles. & Ser superusuario; eludir la auditoría del sistema. \\
        \texttt{app\_core\_rw} & \texttt{LOGIN}, \texttt{NOSUPERUSER}, \texttt{NOCREATEDB}, \texttt{NOCREATEROLE}, \texttt{NOINHERIT} & \texttt{SELECT}, \texttt{INSERT}, \texttt{UPDATE}, \texttt{DELETE}; usar secuencias. & DDL, \texttt{TRUNCATE}, cambios de roles, acceso a funciones administrativas. \\
        \texttt{auditor\_ro} & \texttt{LOGIN}, \texttt{NOSUPERUSER}, \texttt{NOCREATEDB}, \texttt{NOCREATEROLE}, \texttt{NOINHERIT} & \texttt{SELECT} sobre tablas no sensibles y sobre \texttt{vw\_socios\_audit}. & Leer \texttt{socios}, \texttt{cuentas\_ahorro} o \texttt{vw\_cartera\_mora}; modificar datos. \\
        \texttt{migration\_loader} & \texttt{LOGIN}, \texttt{NOSUPERUSER}, \texttt{NOCREATEDB}, \texttt{NOCREATEROLE}, \texttt{NOINHERIT}, \texttt{search\_path='stg'} & Cargar datos en schema \texttt{stg}; leer staging. & Acceso a \texttt{public}; modificar objetos de producción. \\
        \texttt{nicolas}, \texttt{joaquin}, \texttt{david}, \texttt{gustavo} & \texttt{LOGIN}, \texttt{NOSUPERUSER}, \texttt{NOCREATEDB}, \texttt{NOCREATEROLE}, \texttt{NOINHERIT} & \texttt{CONNECT}, \texttt{USAGE} en \texttt{public}; \texttt{ALL} sobre tablas, secuencias y funciones de \texttt{public}. & Ser superusuario; eludir \texttt{pg\_hba}. \\
        \bottomrule
      \end{tabular}
    \end{tablaAPA}%
  }{Elaboración propia con datos de \texttt{evidence/03\_pruebas\_y\_hardening.txt}.}
\end{table}

\capturaAPA{fig:cap-roles}{Privilegios efectivos de los roles operativos del destino}%
  {captura-roles-privilegios.png}%
  {Salida de \texttt{has\_schema\_privilege} y \texttt{has\_table\_privilege} para \texttt{admin\_dba}, \texttt{app\_core\_rw} y \texttt{auditor\_ro}.}

El rol \texttt{admin\_dba} no es superusuario. Puede crear bases y roles,
pero no puede leer archivos del sistema, ejecutar comandos del sistema
operativo ni eludir las reglas de \texttt{pg\_hba}. Esa distinción es
importante: en muchos equipos se usa el superusuario para todo, lo que
convierte cualquier error en un incidente grave. El rol
\texttt{app\_core\_rw} no puede ejecutar \texttt{TRUNCATE} ni DDL, de modo
que una inyección SQL en la aplicación no puede borrar tablas ni alterar el
esquema. El rol \texttt{auditor\_ro} no puede leer las tablas base con datos
personales; solo accede a la vista enmascarada y a las vistas agregadas que no
exponen identificadores completos.

\begin{lstlisting}[language=SQL,caption={Extracto de revocación y creación de roles},label={lst:roles}]
REVOKE ALL ON SCHEMA public FROM PUBLIC;
REVOKE ALL ON ALL TABLES IN SCHEMA public FROM PUBLIC;
REVOKE ALL ON ALL SEQUENCES IN SCHEMA public FROM PUBLIC;
REVOKE ALL ON ALL FUNCTIONS IN SCHEMA public FROM PUBLIC;
REVOKE CONNECT ON DATABASE coop_horizonte_core FROM PUBLIC;

CREATE ROLE admin_dba WITH LOGIN NOSUPERUSER CREATEDB CREATEROLE NOINHERIT
  PASSWORD '...';
CREATE ROLE app_core_rw WITH LOGIN NOSUPERUSER NOCREATEDB NOCREATEROLE
  NOINHERIT PASSWORD '...';
CREATE ROLE auditor_ro WITH LOGIN NOSUPERUSER NOCREATEDB NOCREATEROLE
  NOINHERIT PASSWORD '...';
CREATE ROLE migration_loader WITH LOGIN NOSUPERUSER NOCREATEDB NOCREATEROLE
  NOINHERIT PASSWORD '...';

DO $$
DECLARE
  admin_names text[] := ARRAY['nicolas', 'joaquin', 'david', 'gustavo'];
  admin_name text;
BEGIN
  FOREACH admin_name IN ARRAY admin_names LOOP
    EXECUTE format('CREATE ROLE %I WITH LOGIN NOSUPERUSER NOCREATEDB NOCREATEROLE NOINHERIT PASSWORD %L', admin_name, 'sudoers123');
    EXECUTE format('GRANT CONNECT ON DATABASE coop_horizonte_core TO %I', admin_name);
    EXECUTE format('GRANT USAGE ON SCHEMA public TO %I', admin_name);
    EXECUTE format('GRANT ALL ON ALL TABLES IN SCHEMA public TO %I', admin_name);
    EXECUTE format('GRANT ALL ON ALL SEQUENCES IN SCHEMA public TO %I', admin_name);
    EXECUTE format('GRANT ALL ON ALL FUNCTIONS IN SCHEMA public TO %I', admin_name);
    EXECUTE format('ALTER DEFAULT PRIVILEGES IN SCHEMA public GRANT ALL ON TABLES TO %I', admin_name);
    EXECUTE format('ALTER DEFAULT PRIVILEGES IN SCHEMA public GRANT ALL ON SEQUENCES TO %I', admin_name);
    EXECUTE format('ALTER DEFAULT PRIVILEGES IN SCHEMA public GRANT ALL ON FUNCTIONS TO %I', admin_name);
  END LOOP;
END $$;

GRANT CONNECT ON DATABASE coop_horizonte_core TO admin_dba, app_core_rw, auditor_ro, migration_loader;
GRANT USAGE ON SCHEMA public TO admin_dba, app_core_rw, auditor_ro, migration_loader;

GRANT SELECT, INSERT, UPDATE, DELETE ON ALL TABLES IN SCHEMA public TO app_core_rw;
GRANT USAGE, SELECT ON ALL SEQUENCES IN SCHEMA public TO app_core_rw;

GRANT SELECT ON ALL TABLES IN SCHEMA public TO auditor_ro;
REVOKE SELECT ON socios FROM auditor_ro;
REVOKE SELECT ON cuentas_ahorro FROM auditor_ro;
REVOKE SELECT ON vw_cartera_mora FROM auditor_ro;
GRANT SELECT ON vw_socios_audit TO auditor_ro;
GRANT SELECT ON vw_resumen_saldos_socio TO auditor_ro;
\end{lstlisting}

\subsection{Auditoría y registro de eventos}

El servidor registra conexiones, desconexiones y sentencias DDL. Cada línea
de registro incluye la marca de tiempo, el identificador del proceso, el
usuario y la base de datos, lo que permite correlacionar un evento con una
sesión. Esta
configuración no sustituye a una auditoría completa de DML: para eso se
recomienda \texttt{pgaudit} o una solución de auditoría a nivel de aplicación.
Sin embargo, cubre los eventos más relevantes para el proyecto: quién se
conectó, cuándo y qué cambios de estructura se ejecutaron.

El registro tiene un costo: escribe en disco y puede crecer. En producción se
recomienda rotar los archivos, enviarlos a un recolector central y definir una
retención mínima. La configuración del laboratorio prioriza la evidencia y la
trazabilidad; la operación real debe equilibrar detalle y volumen.

\subsection{Enmascaramiento de datos sensibles}

La vista \texttt{vw\_socios\_audit} permite que el rol \texttt{auditor\_ro}
consulte información de socios sin ver documentos, teléfonos, correos ni
números de cuenta completos. El enmascaramiento se implementa con
\texttt{CASE} y funciones de cadena: se conserva el sufijo para que el
auditor pueda correlacionar registros, pero se ocultan los dígitos
significativos. La Tabla~\ref{tab:enmascaramiento} muestra las reglas.

\begin{table}[H]
  \centering
  \caption{Reglas de enmascaramiento en \texttt{vw\_socios\_audit}}
  \label{tab:enmascaramiento}
  \contenidoConFuente{%
    \begin{tablaAPA}
      \begin{tabular}{@{}L{3.2cm}L{4.2cm}L{7.2cm}@{}}
        \toprule
        \textbf{Campo} & \textbf{Expresión} & \textbf{Ejemplo de salida} \\
        \midrule
        CI/NIT & \texttt{'****' || right(numero\_identificacion, 4)} & \texttt{****6-OR} \\
        Teléfono & \texttt{'****' || right(telefono, 4)} & \texttt{****4877} \\
        Correo & \texttt{'****@' || split\_part(email, '@', 2)} & \texttt{****@horizonte.pymes} \\
        Número de cuenta & \texttt{'****' || right(num\_cuenta, 4)} & \texttt{****9292} \\
        \bottomrule
      \end{tabular}
    \end{tablaAPA}%
  }{Elaboración propia con datos de \texttt{evidence/03\_pruebas\_y\_hardening.txt}.}
\end{table}

\capturaAPA{fig:cap-vista}{Vista enmascarada y acceso denegado a la tabla base}%
  {captura-vista-enmascarada.png}%
  {Consulta a \texttt{vw\_socios\_audit} con valores enmascarados y error \texttt{permission denied} al leer \texttt{socios}.}

La vista agrupa las cuentas por socio con \texttt{string\_agg}, de modo que un
socio con dos cuentas aparece en una sola fila con ambos números enmascarados.
El rol \texttt{auditor\_ro} no tiene \texttt{SELECT} sobre \texttt{socios} ni
sobre \texttt{cuentas\_ahorro}; si intenta consultarlas directamente, recibe
\texttt{permission denied}. Esa separación es la que hace que el
enmascaramiento sea efectivo: no basta con ocultar columnas en una vista si el
usuario puede leer la tabla original.

\subsection{Justificación normativa}

Las medidas se alinean con marcos reconocidos. El \emph{CIS PostgreSQL
Benchmark} recomienda exigir TLS, usar autenticación fuerte, revocar
privilegios de \texttt{PUBLIC}, separar roles y habilitar registros de
conexión \citep{cis2024postgres}. El \emph{OWASP Application Security
Verification Standard} exige control de acceso, cifrado en tránsito y
protección de datos sensibles \citep{owasp2023asvs}. El \emph{NIST SP 800-53}
establece controles de acceso, auditoría y protección de la información
\citep{nist80053}, y el \emph{NIST SP 800-63B} recomienda autenticadores
resistentes a la interceptación \citep{nist80063b}. El proyecto no certifica
cumplimiento, pero adopta estas prácticas como referencia de diseño.

\section{Evidencias de verificación del hardening}
\label{sec:evidencias-hardening}

El script \texttt{scripts/03\_verify\_all.sh} ejecuta 75 comprobaciones y
genera \texttt{evidence/03\_pruebas\_y\_hardening.txt}. El resultado fue
75 \texttt{PASS} y 0 \texttt{FAIL}. La Tabla~\ref{tab:pruebas-hardening}
agrupa las comprobaciones por categoría y resume el resultado. Todas las
consultas se ejecutaron contra el servidor real, no contra una maqueta.

\begin{table}[H]
  \centering
  \caption{Resumen de las pruebas de hardening}
  \label{tab:pruebas-hardening}
  \contenidoConFuente{%
    \begin{tablaAPA}
      \begin{tabular}{@{}L{4.6cm}C{1.8cm}C{1.8cm}L{6.4cm}@{}}
        \toprule
        \textbf{Categoría} & \textbf{PASS} & \textbf{FAIL} & \textbf{Evidencia principal} \\
        \midrule
        Parámetros de seguridad & 5 & 0 & \texttt{ssl=on}, \texttt{scram-sha-256}, logs activos. \\
        Reglas \texttt{pg\_hba} & 4 & 0 & Sin SSL rechazado; SSL y \texttt{verify-ca} aceptados; \texttt{postgres} remoto rechazado. \\
        Roles y privilegios & 27 & 0 & 7 roles auditados; 16 checks de los 4 usuarios admin. \\
        Enmascaramiento & 3 & 0 & Valores con \texttt{****}; sin identificadores completos. \\
        Conteos origen--destino & 8 & 0 & 8 tablas con conteos idénticos. \\
        Integridad referencial & 10 & 0 & 9 claves foráneas; 0 huérfanos. \\
        Staging, loader y secuencias & 18 & 0 & Schema \texttt{stg}; pgloader; \texttt{setval}; \texttt{ANALYZE}. \\
        \midrule
        \textbf{Total} & \textbf{75} & \textbf{0} & Todas las pruebas pasaron. \\
        \bottomrule
      \end{tabular}
    \end{tablaAPA}%
  }{Elaboración propia con datos de \texttt{evidence/03\_pruebas\_y\_hardening.txt}.}
\end{table}

El siguiente fragmento reproduce las salidas más relevantes. Se conserva el
texto original para que el lector pueda contrastarlo con el archivo de
evidencia.

\begin{lstlisting}[style=terminal,caption={Evidencias reales de hardening (extracto)},label={lst:evidencia-hardening}]
== 3. Rechazo de conexiones sin SSL ==
PASS  Conexion TCP sin SSL rechazada
PASS  Conexion TCP con SSL (sslmode=require)
PASS  Conexion TCP con SSL verificado (sslmode=verify-ca)
PASS  Acceso remoto directo de postgres rechazado

== 4. Roles y privilegios de minimo privilegio ==
PASS  admin_dba no es superusuario                               (valor=f)
PASS  app_core_rw no es superusuario                             (valor=f)
PASS  app_core_rw puede iniciar sesion                           (valor=t)
PASS  auditor_ro puede iniciar sesion                            (valor=t)
PASS  PUBLIC sin CREATE en schema public                         (valor=f)
PASS  auditor_ro con USAGE en public                             (valor=t)
PASS  app_core_rw sin CREATE en public                           (valor=f)
PASS  auditor_ro no lee socios                                   (valor=f)
PASS  auditor_ro lee vw_socios_audit                             (valor=t)
PASS  app_core_rw inserta en socios                              (valor=t)
PASS  app_core_rw no trunca socios                               (valor=f)

== 5. Enmascaramiento de datos sensibles ==
****6-OR|****4877, ****9292
****4-OR|****7514
****3-LP|****4895, ****1669
PASS  vw_socios_audit devuelve valores enmascarados
PASS  No se detectaron identificadores completos en la vista
PASS  auditor_ro no puede leer la tabla socios

== 6. Roles de administradores del equipo ==
PASS  nicolas puede iniciar sesion                                 (valor=t)
PASS  joaquin puede iniciar sesion                                 (valor=t)
PASS  david puede iniciar sesion                                   (valor=t)
PASS  gustavo puede iniciar sesion                                 (valor=t)
PASS  nicolas no es superusuario                                   (valor=f)
PASS  joaquin no es superusuario                                   (valor=f)
PASS  david no es superusuario                                     (valor=f)
PASS  gustavo no es superusuario                                   (valor=f)
PASS  nicolas tiene privilegio de admin en public                  (valor=t)
PASS  joaquin tiene privilegio de admin en public                  (valor=t)
PASS  david tiene privilegio de admin en public                    (valor=t)
PASS  gustavo tiene privilegio de admin en public                  (valor=t)
PASS  conexion TLS real para nicolas                               (valor=t)
PASS  conexion TLS real para joaquin                               (valor=t)
PASS  conexion TLS real para david                                 (valor=t)
PASS  conexion TLS real para gustavo                               (valor=t)

== 7. Staging, loader y secuencias ==
PASS  schema stg existe                                            (valor=t)
PASS  migration_loader tiene USAGE en stg                          (valor=t)
PASS  migration_loader sin acceso a public                         (valor=f)
PASS  pgloader cargo tablas sin rechazos                           (valor=t)
PASS  conteos origen vs stg coinciden                              (valor=t)
PASS  precheck valida nulos, ENUM, huérfanos, rangos               (valor=t)
PASS  postcheck EXCEPT ALL bidireccional OK                        (valor=t)
PASS  secuencias reajustadas con setval                            (valor=t)
PASS  ANALYZE ejecutado como postgres                              (valor=t)
PASS  staging eliminado despues de la carga                        (valor=t)
PASS  rollback transaccional confirmado                            (valor=t)

== 8. Resumen final ==
PASS: 75
FAIL: 0
\end{lstlisting}

La interpretación es directa. Si \texttt{ssl} estuviera en \texttt{off}, la
primera prueba fallaría. Si \texttt{app\_core\_rw} pudiera truncar
\texttt{socios}, la prueba de privilegios fallaría. Si la vista devolviera un
documento completo, la expresión regular detectaría siete o más dígitos
consecutivos y la prueba de enmascaramiento fallaría. El script no se limita
a imprimir valores: compara con el valor esperado y cuenta los fallos. Esa
es la diferencia entre una demostración y una verificación. La salida completa
de las 75 pruebas se presenta en el Anexo~C
(Figuras~\ref{fig:anexo-pruebas-1} a~\ref{fig:anexo-pruebas-3}).


% --- CAPÍTULO IV -------------------------------------------------------------
\caratulaCapitulo{CAPÍTULO IV}{DESARROLLO DEL TRABAJO}
% ============================================================================
%  CAPÍTULO III: DESARROLLO DEL TRABAJO (parte 1)
% ============================================================================

\section{Plan de migración}

\subsection{Estrategia seleccionada}

La selección de estrategia dependió de tres variables: el volumen (0.28 MB y
1\,469 filas), la ventana de mantenimiento disponible y la criticidad del
negocio. Se evaluaron tres alternativas. La primera, migración \emph{big bang}
con corte controlado, consiste en congelar escrituras en el origen, cargar los
datos a un schema de staging (\texttt{stg}) con pgloader, validar con
\emph{precheck}, aplicar en una sola transacción con rollback garantizado, y
cambiar la cadena de conexión. La segunda, replicación continua con
\emph{Change Data Capture} (CDC), mantiene ambos motores sincronizados y
reduce el tiempo de inactividad, pero requiere herramientas adicionales,
monitoreo permanente y un esfuerzo de ingeniería desproporcionado para el
volumen actual. La tercera, migración por lotes con doble escritura, exige
modificar la aplicación y convivir con dos fuentes de verdad durante semanas.

La Tabla~\ref{tab:estrategias} compara las alternativas. Para este proyecto se
eligió el corte controlado con staging porque el volumen permite copiar y
verificar en segundos, el respaldo y la restauración están probados, y el
rollback es transaccional: si el \emph{precheck} o el \emph{postcheck} falla,
se conserva el origen y el destino no sufre modificaciones parciales. La
ventana de mantenimiento estimada es de 30 a 60 minutos, muy por debajo de las
4 horas que justificarían CDC.

\begin{table}[H]
  \centering
  \caption{Comparación de estrategias de migración}
  \label{tab:estrategias}
  \contenidoConFuente{%
    \begin{tablaAPA}
      \begin{tabular}{@{}L{3.2cm}L{3.4cm}L{3.4cm}L{4.2cm}@{}}
        \toprule
        \textbf{Estrategia} & \textbf{Ventaja} & \textbf{Desventaja} & \textbf{Decisión} \\
        \midrule
        \emph{Big bang} controlado con staging & Simple, verificable, rollback transaccional; pgloader + precheck/postcheck. & Requiere ventana de inactividad y congelamiento de escrituras. & \textbf{Seleccionada}: el volumen y la criticidad lo permiten. \\
        CDC con replicación & Inactividad mínima; útil para volúmenes grandes. & Complejidad operativa y de monitoreo; herramientas adicionales. & Descartada para esta fase. \\
        Doble escritura & Transición gradual. & Cambios en la aplicación; riesgo de divergencia. & Descartada: exige modificar el código. \\
        \bottomrule
      \end{tabular}
    \end{tablaAPA}%
  }{Elaboración propia.}
\end{table}

\subsection{Fases y cronograma}

El plan se organizó en diez fases. La Tabla~\ref{tab:fases} indica el
objetivo, las actividades principales, el entregable y la duración estimada.
Las fases 1 a 10 se ejecutaron en el laboratorio; las fases de validación se
repitieron hasta obtener 75 pruebas en \texttt{PASS}.

\begin{table}[H]
  \centering
  \caption{Fases del plan de migración}
  \label{tab:fases}
  \contenidoConFuente{%
    \begin{tablaAPA}
      \begin{tabular}{@{}C{1.0cm}L{2.9cm}L{6.2cm}L{3.4cm}C{1.3cm}@{}}
        \toprule
        \textbf{Fase} & \textbf{Nombre} & \textbf{Actividades} & \textbf{Entregable} & \textbf{Duración} \\
        \midrule
        1 & Congelamiento & Detener escrituras en MySQL; validar conectividad origen--destino. & Entorno \emph{healthy} congelado. & 5 min \\
        2 & Respaldo origen & \texttt{mysqldump --single-transaction} del origen. & \texttt{backups/origin\_*.sql.gz}. & 10 min \\
        3 & Staging & \texttt{00\_prepare\_staging.sql} crea schema \texttt{stg}. & Schema \texttt{stg} listo. & 5 min \\
        4 & Carga pgloader & pgloader carga 8 tablas a \texttt{stg} (sin índices, FK, secuencias). & Datos en \texttt{stg}. & 10 min \\
        5 & Conteo staging & Conteos origen vs \texttt{stg}. & Conteos coincidentes. & 5 min \\
        6 & Precheck & \texttt{10\_precheck.sql} valida nulos, ENUM, huérfanos, rangos, fechas cero. Aborta si falla. & Validación superada. & 10 min \\
        7 & Carga transaccional & \texttt{20\_load.sql} + \texttt{30\_postcheck.sql} + \texttt{35\_sequences.sql} en una transacción (\texttt{ROLLBACK} si falla). & Datos en \texttt{public} verificados. & 10 min \\
        8 & Estadísticas & \texttt{40\_analyze.sql} como \texttt{postgres} por socket local. & Planificador actualizado. & 5 min \\
        9 & Limpieza & \texttt{50\_cleanup.sql} elimina schema \texttt{stg}. & Entorno limpio. & 5 min \\
        10 & Respaldo destino y restauración & \texttt{pg\_dump -Fc} + AES-256-CBC; prueba de restauración con 55 socios. & Respaldo cifrado verificado. & 15 min \\
        \bottomrule
      \end{tabular}
    \end{tablaAPA}%
  }{Elaboración propia.}
\end{table}

\subsection{Plan de contingencia y reversión}

El plan de contingencia se diseñó antes de ejecutar la migración. La lógica es
simple: el origen es la fuente de verdad hasta que el destino pasa todas las
pruebas; si algo falla, se conserva el origen y se restaura el destino. La
Tabla~\ref{tab:contingencia} define los disparadores, las acciones y los
responsables. El objetivo de tiempo de recuperación (RTO) es de 30 minutos y
el objetivo de punto de recuperación (RPO) es de 0 minutos para los datos
migrados, porque el respaldo se toma inmediatamente antes del corte.

\begin{table}[H]
  \centering
  \caption{Plan de contingencia y reversión}
  \label{tab:contingencia}
  \contenidoConFuente{%
    \begin{tablaAPA}
      \begin{tabular}{@{}L{3.5cm}L{6.6cm}L{4.6cm}@{}}
        \toprule
        \textbf{Disparador} & \textbf{Acción} & \textbf{Responsable} \\
        \midrule
        Conteos origen vs \texttt{stg} no coinciden & Detener, revisar la carga pgloader, restaurar el destino desde el respaldo y repetir. & DBA / ETL \\
        \texttt{PRECHECK FALLIDO} & Abortar antes de tocar \texttt{public}; corregir la validación fallida (nulos, ENUM, huérfanos, rangos, fechas cero) y repetir desde la fase de carga. & DBA \\
        Error en transacción única & \texttt{ROLLBACK} automático; el destino no sufre modificaciones parciales. Revisar el error en \texttt{20\_load.sql} o \texttt{30\_postcheck.sql} y repetir. & DBA \\
        Clave foránea violada & Identificar el huérfano, corregir la causa en el origen o en la transformación, repetir la carga. & DBA \\
        Falla de conexión TLS & Revisar certificados, \texttt{pg\_hba} y \texttt{sslrootcert}; no continuar sin cifrado. & Seguridad \\
        Error de secuencia & Ejecutar \texttt{setval} con el máximo ID y probar un \texttt{INSERT} nuevo. & DBA \\
        Pérdida de datos & Restaurar el respaldo cifrado, validar 55 socios y 1\,469 filas, repetir la migración. & DBA \\
        Incidente de seguridad & Aislar el contenedor, revocar credenciales, revisar logs y notificar. & Seguridad / TI \\
        \bottomrule
      \end{tabular}
    \end{tablaAPA}%
  }{Elaboración propia.}
\end{table}

Los criterios de aceptación fueron: (a) 100\,\% de las filas migradas y
verificadas; (b) cero huérfanos en las 9 relaciones; (c) 75 pruebas en
\texttt{PASS}; (d) conexiones sin SSL rechazadas; (e) \texttt{auditor\_ro} sin
acceso a tablas sensibles; (f) respaldo cifrado restaurado con 55 socios; y
(g) entorno levantado desde cero con un solo comando. Si cualquiera de estos
criterios falla, la migración no se considera aceptada.

\section{Diseño del esquema destino}

\subsection{Principios de diseño}

El esquema destino conserva los nombres de las tablas y columnas del origen
para reducir el impacto en la aplicación, pero cambia los tipos y añade
restricciones. Se aplicaron cinco principios. Primero, usar tipos nativos de
PostgreSQL en lugar de emular los de MySQL. Segundo, declarar la integridad
referencial en el motor. Tercero, mantener la precisión decimal de los montos.
Cuarto, indexar las columnas usadas en reportes y claves foráneas. Quinto,
exponer los datos sensibles solo a través de vistas controladas.

La Figura~\ref{fig:modelo-datos} muestra el modelo relacional simplificado.
Las relaciones principales son: un socio tiene muchas cuentas y muchas
solicitudes; una solicitud aprobada genera un préstamo; un préstamo tiene
muchas cuotas; y las transacciones referencian un socio, una cuenta o un
préstamo según el tipo de movimiento.

\figuraAPA{fig:modelo-datos}{Modelo relacional simplificado del destino}{%
  \begin{tikzpicture}[
      ent/.style={draw, rounded corners, minimum width=2.7cm,
                  minimum height=0.85cm, align=center, font=\scriptsize,
                  fill=blue!4},
      linea/.style={-Stealth, thick, black!70},
      etiqueta/.style={font=\tiny, fill=white, inner sep=1pt}]
    \node[ent] (suc) at (0,0) {\texttt{sucursales}};
    \node[ent] (soc) at (3.4,0) {\texttt{socios}};
    \node[ent] (cta) at (6.8,0) {\texttt{cuentas\_ahorro}};
    \node[ent] (tx)  at (10.2,0) {\texttt{transacciones}};
    \node[ent] (tip) at (0,-2.6) {\texttt{tipos\_credito}};
    \node[ent] (sol) at (3.4,-2.6) {\texttt{solicitudes\_credito}};
    \node[ent] (pre) at (6.8,-2.6) {\texttt{prestamos}};
    \node[ent] (cuo) at (10.2,-2.6) {\texttt{cuotas\_amortizacion}};

    \draw[linea] (soc) -- node[etiqueta,above]{1:N} (cta);
    \draw[linea] (soc) -- node[etiqueta,right]{1:N} (sol);
    \draw[linea] (tip) -- node[etiqueta,above]{1:N} (sol);
    \draw[linea] (sol) -- node[etiqueta,above]{1:1} (pre);
    \draw[linea] (pre) -- node[etiqueta,above]{1:N} (cuo);
    \draw[linea] (cta) -- node[etiqueta,above]{1:N} (tx);
    \draw[linea] (pre) to[bend left=18] node[etiqueta,right]{1:N} (tx);
    \draw[linea, dashed] (soc) to[bend left=12] node[etiqueta,left]{1:N} (pre);
  \end{tikzpicture}%
}{Elaboración propia con base en \texttt{src/target/01\_schema\_target.sql}.}

\subsection{Tablas, claves y restricciones}

La Tabla~\ref{tab:tablas-destino} describe las 8 tablas, su clave primaria,
sus claves foráneas y las restricciones adicionales. Todas las claves
primarias usan identidades \texttt{BIGINT}. Las claves foráneas tienen
políticas explícitas: \texttt{ON DELETE CASCADE} en cuotas (si se elimina un
préstamo, sus cuotas no tienen sentido) y \texttt{ON DELETE SET NULL} en
transacciones (el historial debe conservarse aunque se elimine la referencia).

\begin{table}[H]
  \centering
  \caption{Tablas y restricciones del esquema destino}
  \label{tab:tablas-destino}
  \contenidoConFuente{%
    \begin{tablaAPA}
      \begin{tabular}{@{}L{3.0cm}L{3.2cm}L{4.0cm}L{4.2cm}@{}}
        \toprule
        \textbf{Tabla} & \textbf{Clave primaria} & \textbf{Claves foráneas} & \textbf{Restricciones} \\
        \midrule
        \texttt{sucursales} & \texttt{id} & --- & \texttt{estatus BOOLEAN}. \\
        \texttt{socios} & \texttt{id} & --- & \texttt{UNIQUE(numero\_identificacion)}; \texttt{CHECK(genero)}; \texttt{CHECK(estado)}. \\
        \texttt{cuentas\_ahorro} & \texttt{id} & \texttt{socio\_id} $\rightarrow$ \texttt{socios} & \texttt{UNIQUE(num\_cuenta)}; saldo y tasa no negativos. \\
        \texttt{tipos\_credito} & \texttt{id} & --- & \texttt{CHECK(estado)}; monto y plazo positivos. \\
        \texttt{solicitudes\_credito} & \texttt{id} & \texttt{socio\_id}, \texttt{tipo\_credito\_id} & \texttt{CHECK(status)}; monto y plazo positivos. \\
        \texttt{prestamos} & \texttt{id} & \texttt{solicitud\_id}, \texttt{socio\_id} & \texttt{UNIQUE(solicitud\_id)}; \texttt{CHECK(estado)}. \\
        \texttt{cuotas\_amortizacion} & \texttt{id} & \texttt{prestamo\_id} & \texttt{UNIQUE(prestamo\_id, numero\_cuota)}; \texttt{CHECK(estatus)}. \\
        \texttt{transacciones} & \texttt{id} & \texttt{socio\_id}, \texttt{cuenta\_ahorro\_id}, \texttt{prestamo\_id} & \texttt{CHECK(tipo)}; monto positivo. \\
        \bottomrule
      \end{tabular}
    \end{tablaAPA}%
  }{Elaboración propia con datos de \texttt{src/target/01\_schema\_target.sql}.}
\end{table}

El fragmento de DDL de la Tabla~\ref{tab:tablas-destino} se implementó en
\texttt{src/target/01\_schema\_target.sql}. Se reproducen las definiciones de
\texttt{socios}, \texttt{prestamos} y \texttt{transacciones} para mostrar el
estilo.

\begin{lstlisting}[language=SQL,caption={Definiciones representativas del esquema destino},label={lst:ddl-destino}]
CREATE TABLE socios (
  id BIGINT GENERATED ALWAYS AS IDENTITY PRIMARY KEY,
  nombre_completo VARCHAR(150) NOT NULL,
  numero_identificacion VARCHAR(20) UNIQUE,
  fecha_nacimiento DATE,
  genero VARCHAR(1) NOT NULL DEFAULT 'M'
    CHECK (genero IN ('M', 'F', 'O')),
  estado VARCHAR(10) NOT NULL DEFAULT 'activo'
    CHECK (estado IN ('activo', 'inactivo', 'suspendido')),
  fecha_ingreso DATE NOT NULL,
  telefono VARCHAR(20),
  email VARCHAR(100),
  direccion_residencia TEXT,
  fecha_ultima_actualizacion TIMESTAMPTZ NOT NULL DEFAULT now(),
  estatus BOOLEAN NOT NULL DEFAULT TRUE
);

CREATE TABLE prestamos (
  id BIGINT GENERATED ALWAYS AS IDENTITY PRIMARY KEY,
  solicitud_id BIGINT NOT NULL UNIQUE REFERENCES solicitudes_credito(id),
  socio_id BIGINT NOT NULL REFERENCES socios(id),
  monto_prestado NUMERIC(12,2) NOT NULL CHECK (monto_prestado > 0),
  tasa_interes_anual NUMERIC(5,2) NOT NULL CHECK (tasa_interes_anual >= 0),
  fecha_desembolso TIMESTAMPTZ NOT NULL,
  fecha_fin_plazo DATE NOT NULL,
  estado VARCHAR(10) NOT NULL DEFAULT 'activo'
    CHECK (estado IN ('activo', 'pagado', 'moroso', 'cancelado')),
  total_cuotas BIGINT NOT NULL CHECK (total_cuotas > 0),
  pagado_total BOOLEAN NOT NULL DEFAULT FALSE
);

CREATE TABLE transacciones (
  id BIGINT GENERATED ALWAYS AS IDENTITY PRIMARY KEY,
  tipo VARCHAR(15) NOT NULL
    CHECK (tipo IN ('deposito', 'retiro', 'pago_cuota', 'desembolso', 'transferencia')),
  socio_id BIGINT REFERENCES socios(id) ON DELETE SET NULL,
  cuenta_ahorro_id BIGINT REFERENCES cuentas_ahorro(id) ON DELETE SET NULL,
  prestamo_id BIGINT REFERENCES prestamos(id) ON DELETE SET NULL,
  monto NUMERIC(12,2) NOT NULL CHECK (monto > 0),
  fecha_transaccion TIMESTAMPTZ NOT NULL,
  descripcion VARCHAR(255)
);
\end{lstlisting}

\subsection{Índices}

Se crearon 10 índices no únicos además de las claves primarias y las
restricciones únicas. La selección se basó en las consultas de los reportes y
en las columnas de join. La Tabla~\ref{tab:indices-destino} justifica cada
índice. No se indexó todo: cada índice adicional consume espacio y ralentiza
las escrituras. Se priorizaron las columnas con mayor selectividad y uso.

\begin{table}[H]
  \centering
  \caption{Índices del esquema destino y su justificación}
  \label{tab:indices-destino}
  \contenidoConFuente{%
    \begin{tablaAPA}
      \begin{tabular}{@{}L{4.2cm}L{3.2cm}L{7.0cm}@{}}
        \toprule
        \textbf{Índice} & \textbf{Tabla} & \textbf{Uso previsto} \\
        \midrule
        \texttt{idx\_socio\_estado} & \texttt{socios} & Filtrar socios activos, inactivos o suspendidos. \\
        \texttt{idx\_cuenta\_ahorro\_socio} & \texttt{cuentas\_ahorro} & Reporte de saldos por socio. \\
        \texttt{idx\_solicitud\_tipo} & \texttt{solicitudes\_credito} & Análisis por producto de crédito. \\
        \texttt{idx\_solicitud\_socio} & \texttt{solicitudes\_credito} & Historial de solicitudes por socio. \\
        \texttt{idx\_transaccion\_fecha} & \texttt{transacciones} & Reportes por rango de fechas. \\
        \texttt{idx\_transaccion\_socio} & \texttt{transacciones} & Movimientos por socio. \\
        \texttt{idx\_prestamo\_estado} & \texttt{prestamos} & Cartera por estado (incluye morosos). \\
        \texttt{idx\_prestamo\_socio} & \texttt{prestamos} & Préstamos por socio. \\
        \texttt{idx\_cuota\_prestamo} & \texttt{cuotas\_amortizacion} & Plan de pagos por préstamo. \\
        \texttt{idx\_cuota\_estatus} & \texttt{cuotas\_amortizacion} & Cuotas pendientes, pagadas o atrasadas. \\
        \bottomrule
      \end{tabular}
    \end{tablaAPA}%
  }{Elaboración propia con datos de \texttt{src/target/01\_schema\_target.sql}.}
\end{table}

\subsection{Vistas}

Se recrearon las dos vistas del origen y se añadió la vista de auditoría. La
vista \texttt{vw\_cartera\_mora} conserva la lógica original: selecciona
préstamos en estado \texttt{moroso}, cuenta sus cuotas y suma las atrasadas.
La vista \texttt{vw\_resumen\_saldos\_socio} agrupa cuentas y préstamos por
socio. La vista \texttt{vw\_socios\_audit} es nueva y aplica enmascaramiento.
La Tabla~\ref{tab:vistas-destino} resume las tres.

\begin{table}[H]
  \centering
  \caption{Vistas del esquema destino}
  \label{tab:vistas-destino}
  \contenidoConFuente{%
    \begin{tablaAPA}
      \begin{tabular}{@{}L{4.4cm}L{2.4cm}L{7.4cm}@{}}
        \toprule
        \textbf{Vista} & \textbf{Tipo} & \textbf{Propósito} \\
        \midrule
        \texttt{vw\_cartera\_mora} & Reporte & Préstamos morosos, cuotas atrasadas y monto en mora. \\
        \texttt{vw\_resumen\_saldos\_socio} & Reporte & Saldos activos y cerrados, cantidad de cuentas y préstamos. \\
        \texttt{vw\_socios\_audit} & Auditoría & Datos de socio con CI/NIT, teléfono, correo y cuentas enmascarados. \\
        \bottomrule
      \end{tabular}
    \end{tablaAPA}%
  }{Elaboración propia con datos de \texttt{src/target/01\_schema\_target.sql} y \texttt{02\_hardening\_roles.sql}.}
\end{table}

\section{Implementación del ETL}

\subsection{Arquitectura del proceso}

El ETL se implementó en Python 3 con \texttt{pymysql} y \texttt{psycopg2},
dentro del contenedor \texttt{coop\_horizonte\_migrador}, construido sobre la
imagen \texttt{postgres:16}. El proceso abre una transacción de lectura en
MySQL con \texttt{REPEATABLE READ} y \texttt{START TRANSACTION WITH CONSISTENT
SNAPSHOT}, de modo que todas las tablas se leen desde el mismo punto en el
tiempo. Luego trunca las tablas destino, carga en orden jerárquico, reajusta
las secuencias y compara conteos. La Figura~\ref{fig:flujo-etl} resume el
flujo.

\figuraAPA{fig:flujo-etl}{Flujo del ETL y de la verificación}{%
  \begin{tikzpicture}[
      paso/.style={draw, rounded corners, minimum width=2.7cm,
                   minimum height=0.9cm, align=center, font=\scriptsize,
                   fill=green!4},
      linea/.style={-Stealth, thick}]
    \node[paso] (ext) {Extraer\\snapshot};
    \node[paso, right=0.7cm of ext] (tr) {Transformar\\tipos};
    \node[paso, right=0.7cm of tr] (load) {Cargar\\destino};
    \node[paso, right=0.7cm of load] (seq) {Reajustar\\secuencias};
    \node[paso, right=0.7cm of seq] (ver) {Verificar\\conteos};
    \draw[linea] (ext) -- (tr);
    \draw[linea] (tr) -- (load);
    \draw[linea] (load) -- (seq);
    \draw[linea] (seq) -- (ver);
  \end{tikzpicture}%
}{Elaboración propia con base en \texttt{src/migration/migrate.py}.}

\subsection{Transformaciones aplicadas}

La Tabla~\ref{tab:transformaciones} documenta las transformaciones del ETL.
Cada una responde a una incompatibilidad identificada en el diagnóstico. La
conversión de booleanos usa \texttt{bool(int(valor))}; la de fechas asigna
\texttt{ZoneInfo("America/La\_Paz")}; la de identificadores conserva el valor
original y lo inserta con \texttt{OVERRIDING SYSTEM VALUE}. Los valores
\texttt{NULL} se preservan en todas las columnas.

\begin{table}[H]
  \centering
  \caption{Transformaciones del ETL}
  \label{tab:transformaciones}
  \contenidoConFuente{%
    \begin{tablaAPA}
      \begin{tabular}{@{}L{3.0cm}L{4.4cm}L{6.8cm}@{}}
        \toprule
        \textbf{Columna origen} & \textbf{Transformación} & \textbf{Justificación} \\
        \midrule
        \texttt{estatus}, \texttt{pagado\_total} & \texttt{bool(int(valor))} & \texttt{TINYINT(1)} $\rightarrow$ \texttt{BOOLEAN}. \\
        \texttt{fecha\_*}, \texttt{fecha\_transaccion} & \texttt{valor.replace(tzinfo=America/La\_Paz)} & \texttt{DATETIME} $\rightarrow$ \texttt{TIMESTAMPTZ}. \\
        \texttt{id} & \texttt{OVERRIDING SYSTEM VALUE} & Conservar referencias y evitar colisiones. \\
        \texttt{estado}, \texttt{status}, \texttt{tipo} & Texto sin cambios; validación por \texttt{CHECK} & \texttt{ENUM} $\rightarrow$ \texttt{VARCHAR} + \texttt{CHECK}. \\
        \texttt{monto}, \texttt{saldo}, \texttt{tasa} & \texttt{Decimal} sin redondeo adicional & Mantener precisión \texttt{NUMERIC(12,2)}. \\
        \texttt{NULL} & Se conserva & No inventar valores ni aplicar defaults. \\
        \bottomrule
      \end{tabular}
    \end{tablaAPA}%
  }{Elaboración propia con base en \texttt{src/migration/migrate.py}.}
\end{table}

\subsection{Carga, idempotencia y verificación}

La carga se hace tabla por tabla, en orden de dependencias. Antes de insertar,
el ETL ejecuta \texttt{TRUNCATE} sobre las 8 tablas con \texttt{CASCADE}. No
se usa \texttt{RESTART IDENTITY} porque \texttt{admin\_dba} no es propietario
de las secuencias creadas por las identidades; en su lugar, al final se
ejecuta \texttt{setval} con el máximo ID de cada tabla. Esa decisión surgió de
un error real documentado en la bitácora
(Sección~\ref{sec:bitacora}). La carga es idempotente: se puede repetir sin
duplicar filas porque el \texttt{TRUNCATE} limpia el destino y el snapshot del
origen es consistente.

\begin{lstlisting}[language=Python,caption={Núcleo del ETL: transformación, carga y reajuste de secuencias},label={lst:etl}]
for table, columns, datetime_cols, bool_cols in TABLE_SPECS:
    mysql_cur.execute(f"SELECT {', '.join(columns)} FROM {table} ORDER BY id")
    raw_rows = mysql_cur.fetchall()
    rows = []
    for raw in raw_rows:
        rows.append(tuple(
            transform(value, col in bool_cols, col in datetime_cols)
            for col, value in zip(columns, raw)
        ))
    if rows:
        insert_sql = (
            f"INSERT INTO {table} ({', '.join(columns)}) "
            f"OVERRIDING SYSTEM VALUE VALUES %s"
        )
        execute_values(pg_cur, insert_sql, rows, page_size=500)

for table, _, _, _ in TABLE_SPECS:
    pg_cur.execute(
        f"SELECT setval(pg_get_serial_sequence(%s, 'id'), "
        f"COALESCE(MAX(id), 1), MAX(id) IS NOT NULL) FROM {table}",
        (table,),
    )
\end{lstlisting}

La verificación compara \texttt{COUNT(*)} en cada tabla del origen y del
destino. Si alguna diferencia existe, el ETL hace \texttt{rollback} y devuelve
código de salida 1. En la ejecución real, las 8 tablas coincidieron y el total
fue 1\,469 filas. La Tabla~\ref{tab:ejecucion-etl} muestra los tiempos por
tabla. La tabla más costosa fue \texttt{transacciones} (740 filas, 0.022 s),
seguida de \texttt{cuotas\_amortizacion} (516 filas, 0.014 s). El tiempo total
de carga fue inferior a 0.06 s, sin contar la conexión y la verificación.

\begin{table}[H]
  \centering
  \caption{Tiempos de carga del ETL por tabla}
  \label{tab:ejecucion-etl}
  \contenidoConFuente{%
    \begin{tablaAPA}
      \begin{tabular}{@{}L{6.4cm}R{2.4cm}R{3.0cm}@{}}
        \toprule
        \textbf{Tabla} & \textbf{Filas} & \textbf{Tiempo (s)} \\
        \midrule
        \texttt{tipos\_credito} & 5 & 0.003 \\
        \texttt{sucursales} & 5 & 0.002 \\
        \texttt{socios} & 55 & 0.003 \\
        \texttt{cuentas\_ahorro} & 66 & 0.003 \\
        \texttt{solicitudes\_credito} & 50 & 0.003 \\
        \texttt{prestamos} & 32 & 0.002 \\
        \texttt{cuotas\_amortizacion} & 516 & 0.014 \\
        \texttt{transacciones} & 740 & 0.022 \\
        \midrule
        \textbf{Total} & \textbf{1\,469} & \textbf{0.052} \\
        \bottomrule
      \end{tabular}
    \end{tablaAPA}%
  }{Elaboración propia con datos de \texttt{evidence/02\_ejecucion\_migracion.txt}.}
\end{table}

\section{Orquestación Docker y automatización}

\subsection{Servicios y aislamiento}

El archivo \texttt{docker-compose.yml} define tres servicios principales y un
perfil de documentación. La Tabla~\ref{tab:servicios} describe cada uno. Los
puertos se publican solo en \texttt{127.0.0.1}, de modo que la base no queda
expuesta a la red local. Los datos persisten en volúmenes nombrados y los
archivos de inicialización se montan como solo lectura.

\begin{table}[H]
  \centering
  \caption{Servicios de Docker Compose}
  \label{tab:servicios}
  \contenidoConFuente{%
    \begin{tablaAPA}
      \begin{tabular}{@{}L{2.8cm}L{3.0cm}L{4.2cm}L{4.4cm}@{}}
        \toprule
        \textbf{Servicio} & \textbf{Imagen} & \textbf{Función} & \textbf{Puerto / volumen} \\
        \midrule
        \texttt{db-origen} & \texttt{mysql:5.7} & Motor legacy con esquema y datos. & \texttt{127.0.0.1:3306}; \texttt{mysql\_data}. \\
        \texttt{db-destino} & \texttt{postgres:16} & Motor destino endurecido. & \texttt{127.0.0.1:5432}; \texttt{postgres\_data}. \\
        \texttt{migrador} & \texttt{postgres:16} + Python & ETL, respaldo y pruebas. & Sin puerto; red interna. \\
        \texttt{docs} & \texttt{texlive/texlive} & Compilación del informe. & Perfil \texttt{docs}. \\
        \bottomrule
      \end{tabular}
    \end{tablaAPA}%
  }{Elaboración propia con datos de \texttt{docker-compose.yml}.}
\end{table}

Los servicios comparten la red \texttt{red\_cooperativa}. El migrador resuelve
los nombres \texttt{db-origen} y \texttt{db-destino} por DNS interno de
Docker, lo que evita direcciones IP fijas. El destino monta
\texttt{postgresql.conf}, \texttt{pg\_hba.conf}, los certificados y los scripts
de inicialización. El entrypoint copia los certificados a una ruta escribible,
ajusta propietario y permisos, y delega en el entrypoint oficial de
PostgreSQL. Esa pequeña capa fue necesaria porque el volumen de datos no
permite escribir directamente en \texttt{/etc/postgresql}.

\subsection{Scripts de automatización}

La Tabla~\ref{tab:scripts} describe los cuatro scripts de soporte. Cada uno
tiene una responsabilidad única y guarda su salida en \texttt{evidence/}. El
script \texttt{lib.sh} centraliza funciones comunes: carga del \texttt{.env},
espera de servicios \emph{healthy}, ejecución de \texttt{mysql} y
\texttt{psql}, y conteos. Así se evita duplicar lógica y se facilita el
mantenimiento.

\begin{table}[H]
  \centering
  \caption{Scripts de automatización}
  \label{tab:scripts}
  \contenidoConFuente{%
    \begin{tablaAPA}
      \begin{tabular}{@{}L{4.4cm}L{9.6cm}@{}}
        \toprule
        \textbf{Script} & \textbf{Responsabilidad} \\
        \midrule
        \texttt{scripts/lib.sh} & Funciones comunes: entorno, espera, consultas y conteos. \\
        \texttt{scripts/01\_audit\_origin.sh} & Inventario, tamaño, objetos, usuarios, privilegios e incompatibilidades. \\
        \texttt{scripts/02\_migrate\_and\_backup.sh} & Respaldo del origen, pgloader + SQL, respaldo cifrado y prueba de restauración. \\
        \texttt{scripts/03\_verify\_all.sh} & 75 pruebas de hardening, roles, enmascaramiento, conteos, integridad, staging y secuencias. \\
        \bottomrule
      \end{tabular}
    \end{tablaAPA}%
  }{Elaboración propia.}
\end{table}

\subsection{Reproducibilidad}

El entorno se levanta con \texttt{docker compose up -d --build}. La primera
vez, MySQL ejecuta \texttt{01\_schema\_origin.sql} y
\texttt{02\_seed\_origin.sql}; PostgreSQL ejecuta
\texttt{01\_schema\_target.sql} y \texttt{02\_hardening\_roles.sql}. Los
healthchecks esperan a que ambos motores estén listos antes de considerar el
servicio disponible. Para repetir la migración desde cero se ejecuta
\texttt{docker compose down -v} y se vuelve a levantar; los volúmenes se
recrean y los datos se generan de nuevo con la misma semilla. Esa
reproducibilidad es la base de la confianza: cualquier integrante puede
verificar los mismos resultados sin depender de una máquina específica.

% ============================================================================
%  CAPÍTULO III: DESARROLLO DEL TRABAJO (parte 2)
% ============================================================================

\section{Ejecución de la migración}

\subsection{Conteos previos y respaldo del origen}

La ejecución comenzó con la verificación de los servicios y el conteo de las
8 tablas del origen. El script \texttt{02\_migrate\_and\_backup.sh} esperó a
que \texttt{db-origen} y \texttt{db-destino} estuvieran \emph{healthy} y a que
el migrador estuviera en ejecución. Luego generó un respaldo del origen con
\texttt{mysqldump --single-transaction --routines --triggers --events},
comprimido con \texttt{gzip -9}. El archivo resultante fue
\texttt{backups/origin\_20261120\_080940.sql.gz}, de 32 KB, con SHA-256
\texttt{6a2f3a958c4732743302299eb7e4f5d722924f749dbf6fa9764a2936428be3b3}.
La Tabla~\ref{tab:conteos-previos} reproduce los conteos.

\begin{table}[H]
  \centering
  \caption{Conteos del origen antes de la migración}
  \label{tab:conteos-previos}
  \contenidoConFuente{%
    \begin{tablaAPA}
      \begin{tabular}{@{}L{6.4cm}R{3.0cm}@{}}
        \toprule
        \textbf{Tabla} & \textbf{Filas} \\
        \midrule
        \texttt{sucursales} & 5 \\
        \texttt{socios} & 55 \\
        \texttt{cuentas\_ahorro} & 66 \\
        \texttt{tipos\_credito} & 5 \\
        \texttt{solicitudes\_credito} & 50 \\
        \texttt{prestamos} & 32 \\
        \texttt{cuotas\_amortizacion} & 516 \\
        \texttt{transacciones} & 740 \\
        \midrule
        \textbf{Total} & \textbf{1\,469} \\
        \bottomrule
      \end{tabular}
    \end{tablaAPA}%
  }{Elaboración propia con datos de \texttt{evidence/02\_ejecucion\_migracion.txt}.}
\end{table}

\subsection{Ejecución del ETL}

La migración se ejecutó mediante \texttt{scripts/02\_migrate\_and\_backup.sh},
que orquesta pgloader con schema de staging, precheck, carga transaccional,
postcheck, reajuste de secuencias, \texttt{ANALYZE} y limpieza. La salida real
se reproduce a continuación. Se observa el congelamiento de escrituras, el
respaldo del origen, la carga pgloader a \texttt{stg}, la validación de
precheck, la carga en una sola transacción, el reajuste de secuencias y el
respaldo cifrado del destino. El total migrado fue 1\,469 filas con 0 rechazos
y 0 errores.

\begin{lstlisting}[style=terminal,caption={Salida real de la migración con pgloader + staging},label={lst:salida-etl}]
== 1. Congelamiento de escrituras ==
Congelamiento activado en MySQL (coop_horizonte_legacy)

== 2. Respaldo del origen ==
mysqldump --single-transaction --routines --triggers --events
Archivo: backups/origin_20261120_080940.sql.gz (32 KB)
SHA-256: 6a2f3a958c4732743302299eb7e4f5d722924f749dbf6fa9764a2936428be3b3

== 3. Preparación de staging ==
psql -v ON_ERROR_STOP=1 -f sql/00_prepare_staging.sql
CREATE SCHEMA stg
ALTER DEFAULT PRIVILEGES IN SCHEMA stg GRANT ALL ON TABLES TO migration_loader

== 4. Carga pgloader a staging ==
pgloader mysql_to_stage.load
... connecting to MySQL ...
... connecting to PostgreSQL ...
types_credito:       5 rows
sucursales:          5 rows
socios:             55 rows
cuentas_ahorro:     66 rows
solicitudes_credito: 50 rows
prestamos:          32 rows
cuotas_amortizacion: 516 rows
transacciones:      740 rows
Total:            1469 rows, 0 rejects, 0 errors
ALTER SCHEMA 'coop_horizonte_legacy' RENAME TO 'stg'

== 5. Conteos origen vs staging ==
tipos_credito:       5  = 5  OK
sucursales:          5  = 5  OK
socios:             55  = 55  OK
cuentas_ahorro:     66  = 66  OK
solicitudes_credito: 50  = 50  OK
prestamos:          32  = 32  OK
cuotas_amortizacion: 516 = 516 OK
transacciones:      740 = 740 OK

== 6. Precheck ==
psql -v ON_ERROR_STOP=1 -f sql/10_precheck.sql
PRECHECK OK: nulos validados
PRECHECK OK: ENUM dentro de dominio
PRECHECK OK: sin huérfanos en claves foráneas
PRECHECK OK: rangos y fechas cero validados

== 7. Carga transaccional (20_load + 30_postcheck + 35_sequences) ==
psql --single-transaction -v ON_ERROR_STOP=1 -f sql/20_load.sql
psql --single-transaction -v ON_ERROR_STOP=1 -f sql/30_postcheck.sql
POSTCHECK OK: EXCEPT ALL bidireccional sin diferencias
psql --single-transaction -v ON_ERROR_STOP=1 -f sql/35_sequences.sql
setval OK: 8 secuencias reajustadas

== 8. Estadísticas ==
psql -U postgres -f sql/40_analyze.sql
ANALYZE OK

== 9. Limpieza de staging ==
psql -v ON_ERROR_STOP=1 -f sql/50_cleanup.sql
DROP SCHEMA stg CASCADE

== 10. Respaldo cifrado del destino ==
pg_dump -Fc coop_horizonte_core | openssl enc -aes-256-cbc -pbkdf2 -iter 200000 -salt
Archivo: backups/target_20261120_080940.dump.enc (68 KB)
SHA-256: 7017f6905da2e9b317eeb57885cb9d5ef60dff4a90c6efa6de160bf44c1c4d22

== 11. Prueba de restauración ==
Descifrado OK: 68K
pg_restore en coop_horizonte_restore
Filas de socios restauradas: 55
Restauracion: OK (coincide con el origen)

== Total migrado ==
1469 filas, 0 rechazos, 0 errores
\end{lstlisting}

\capturaAPA{fig:cap-etl}{Ejecución del ETL hacia PostgreSQL 16}%
  {captura-ejecucion-etl.png}%
  {Salida de \texttt{scripts/02\_migrate\_and\_backup.sh} con la carga de las 8 tablas y la verificación de conteos.}

La salida completa del ETL se presenta en el Anexo~C
(Figura~\ref{fig:anexo-migracion-1}).

\subsection{Respaldo cifrado del destino}

Inmediatamente después de la carga, el script generó un respaldo del destino
con \texttt{pg\_dump -Fc} y lo cifró con OpenSSL en modo AES-256-CBC, PBKDF2
con 200\,000 iteraciones y \emph{salt}. El archivo resultante fue
\texttt{backups/target\_20261120\_080940.dump.enc}, de 68 KB, con SHA-256
\texttt{7017f6905da2e9b317eeb57885cb9d5ef60dff4a90c6efa6de160bf44c1c4d22}.
La Tabla~\ref{tab:respaldos} resume los dos respaldos.

\begin{table}[H]
  \centering
  \caption{Respaldos generados durante la migración}
  \label{tab:respaldos}
  \contenidoConFuente{%
    \begin{tablaAPA}
      \begin{tabular}{@{}L{3.4cm}L{2.0cm}L{2.0cm}L{5.4cm}@{}}
        \toprule
        \textbf{Archivo} & \textbf{Tamaño} & \textbf{Formato} & \textbf{SHA-256 (prefijo)} \\
        \midrule
        \texttt{origin\_20261120\_080940.sql.gz} & 32 KB & \texttt{mysqldump} + gzip & \texttt{6a2f3a95\ldots} \\
        \texttt{target\_20261120\_080940.dump.enc} & 68 KB & \texttt{pg\_dump -Fc} + AES-256-CBC & \texttt{7017f690\ldots} \\
        \bottomrule
      \end{tabular}
    \end{tablaAPA}%
  }{Elaboración propia con datos de \texttt{evidence/02\_ejecucion\_migracion.txt}.}
\end{table}

\subsection{Prueba de restauración}

La prueba de restauración descifró el respaldo, creó la base temporal
\texttt{coop\_horizonte\_restore}, ejecutó \texttt{pg\_restore} y consultó el
conteo de socios. El resultado fue 55, idéntico al origen. La base temporal se
eliminó al finalizar. Esta prueba es fundamental: un respaldo que no se
restaura no es una garantía, es una suposición. El procedimiento se puede
repetir con cualquier respaldo archivado.

\begin{lstlisting}[style=terminal,caption={Prueba real de restauración del respaldo cifrado},label={lst:restauracion}]
== Prueba de restauracion del respaldo cifrado ==
Descifrado OK: 68K
NOTICE:  database "coop_horizonte_restore" does not exist, skipping
DROP DATABASE
CREATE DATABASE
Filas de socios restauradas: 55
Restauracion: OK (coincide con el origen)
\end{lstlisting}

\capturaAPA{fig:cap-restauracion}{Prueba de restauración del respaldo cifrado}%
  {captura-restauracion.png}%
  {Descifrado del respaldo, \texttt{pg\_restore} en la base temporal y 55 filas de socios restauradas.}

El respaldo cifrado y la restauración completa se presentan en el Anexo~C
(Figura~\ref{fig:anexo-migracion-2}).

\section{Respaldo cifrado y recuperación}

\subsection{Estrategia de respaldo}

La estrategia combina dos formatos. El origen se respalda con
\texttt{mysqldump}, que produce un archivo SQL lógico, portable y fácil de
inspeccionar; se comprime con gzip para reducir espacio. El destino se
respalda con \texttt{pg\_dump -Fc}, el formato personalizado de PostgreSQL,
que permite restaurar selectivamente y comprimir internamente. El respaldo
del destino se cifra con AES-256-CBC porque contiene datos personales y
financieros. La Tabla~\ref{tab:respaldo-controles} relaciona cada control con
su propósito.

\begin{table}[H]
  \centering
  \caption{Controles del respaldo y la recuperación}
  \label{tab:respaldo-controles}
  \contenidoConFuente{%
    \begin{tablaAPA}
      \begin{tabular}{@{}L{4.0cm}L{5.2cm}L{5.0cm}@{}}
        \toprule
        \textbf{Control} & \textbf{Implementación} & \textbf{Propósito} \\
        \midrule
        Consistencia & \texttt{mysqldump --single-transaction} & Respaldo sin bloquear la operación. \\
        Compresión & \texttt{gzip -9} en el origen; \texttt{-Fc} en el destino & Reducir tamaño y facilitar transferencia. \\
        Cifrado & AES-256-CBC, PBKDF2 200\,000 iteraciones, \emph{salt} & Proteger la confidencialidad en reposo. \\
        Integridad & SHA-256 de cada archivo & Detectar corrupción o alteración. \\
        Restauración & \texttt{pg\_restore} en base temporal y conteo & Verificar que el respaldo es utilizable. \\
        Retención & Política propuesta de 30 días diarios y 12 meses mensuales & Cumplir operación y auditoría. \\
        \bottomrule
      \end{tabular}
    \end{tablaAPA}%
  }{Elaboración propia.}
\end{table}

\subsection{Procedimiento de recuperación}

El procedimiento de recuperación tiene cinco pasos: (1) identificar el
respaldo y verificar su SHA-256; (2) descifrar con la frase de paso correcta;
(3) crear una base vacía; (4) restaurar con \texttt{pg\_restore}; y (5)
validar conteos y claves foráneas antes de abrir el servicio. En producción,
el paso 5 debe incluir una consulta de negocio (por ejemplo, saldos por
socio) y una revisión de los logs. El RTO estimado es de 30 minutos para el
volumen actual; el RPO es de 0 minutos para los datos migrados porque el
respaldo se toma en el corte.

\subsection{Gestión de claves}

En el laboratorio, la frase de paso se guarda en la variable
\texttt{BACKUP\_PASSPHRASE} del archivo \texttt{.env}, que no se versiona.
En producción, esa práctica no es suficiente: se recomienda usar un gestor de
secretos (Vault, AWS Secrets Manager, Azure Key Vault o equivalente), rotar la
frase de paso periódicamente y separar la custodia de la clave de la custodia
del respaldo. El cifrado simétrico protege el archivo, pero si la clave y el
archivo están en el mismo servidor comprometido, la protección se reduce. La
regla 3-2-1 (tres copias, dos medios, una fuera del sitio) sigue siendo una
buena referencia.

\section{Pruebas en destino}

\subsection{Resultado global}

El script \texttt{03\_verify\_all.sh} ejecutó 75 comprobaciones y obtuvo
75 \texttt{PASS} y 0 \texttt{FAIL}. La Tabla~\ref{tab:pruebas-global} resume
los grupos. Las pruebas no se limitan a verificar que el servidor responde:
comprueban el valor efectivo de los parámetros, el rechazo de conexiones sin
SSL, los privilegios de cada rol, el enmascaramiento, los conteos, las claves
foráneas, el staging, el loader y las secuencias, y el tiempo de ejecución de
las vistas.

\begin{table}[H]
  \centering
  \caption{Resumen de las pruebas en destino}
  \label{tab:pruebas-global}
  \contenidoConFuente{%
    \begin{tablaAPA}
      \begin{tabular}{@{}L{5.6cm}C{2.0cm}C{2.0cm}L{4.4cm}@{}}
        \toprule
        \textbf{Grupo} & \textbf{Pruebas} & \textbf{PASS} & \textbf{Resultado} \\
        \midrule
        Parámetros de seguridad & 5 & 5 & \texttt{ssl}, \texttt{scram}, logs. \\
        \texttt{pg\_hba} y conexiones & 4 & 4 & Sin SSL rechazado; \texttt{verify-ca} OK. \\
        Roles y privilegios & 27 & 27 & 7 roles auditados (3 operación + 1 loader + 4 admin). \\
        Enmascaramiento & 3 & 3 & CI/NIT y cuentas ocultas. \\
        Conteos & 8 & 8 & 8 tablas idénticas. \\
        Integridad referencial & 10 & 10 & 9 FK, 0 huérfanos. \\
        Staging, loader y secuencias & 18 & 18 & Schema \texttt{stg}; pgloader; \texttt{setval}; \texttt{ANALYZE}. \\
        \midrule
        \textbf{Total} & \textbf{75} & \textbf{75} & \textbf{0 FAIL}. \\
        \bottomrule
      \end{tabular}
    \end{tablaAPA}%
  }{Elaboración propia con datos de \texttt{evidence/03\_pruebas\_y\_hardening.txt}.}
\end{table}

\capturaAPA{fig:cap-pruebas}{Resumen de la verificación en destino}%
  {captura-pruebas-resumen.png}%
  {Salida final de \texttt{scripts/03\_verify\_all.sh}: \texttt{PASS: 75} y \texttt{FAIL: 0}.}

Las capturas completas de las pruebas se presentan en el Anexo~C
(Figuras~\ref{fig:anexo-pruebas-1} a~\ref{fig:anexo-pruebas-3}).

\subsection{Conteos origen--destino}

La Tabla~\ref{tab:comparativa-conteos} compara los conteos. Las 8 filas
coinciden exactamente. La coincidencia no demuestra por sí sola que los
valores sean correctos, pero junto con la verificación de claves foráneas y
la restauración del respaldo proporciona una evidencia sólida de que la carga
fue completa.

\begin{table}[H]
  \centering
  \caption{Comparativa de conteos origen vs. destino}
  \label{tab:comparativa-conteos}
  \contenidoConFuente{%
    \begin{tablaAPA}
      \begin{tabular}{@{}L{6.0cm}R{2.4cm}R{2.4cm}C{2.2cm}@{}}
        \toprule
        \textbf{Tabla} & \textbf{Origen} & \textbf{Destino} & \textbf{Estado} \\
        \midrule
        \texttt{sucursales} & 5 & 5 & OK \\
        \texttt{socios} & 55 & 55 & OK \\
        \texttt{cuentas\_ahorro} & 66 & 66 & OK \\
        \texttt{tipos\_credito} & 5 & 5 & OK \\
        \texttt{solicitudes\_credito} & 50 & 50 & OK \\
        \texttt{prestamos} & 32 & 32 & OK \\
        \texttt{cuotas\_amortizacion} & 516 & 516 & OK \\
        \texttt{transacciones} & 740 & 740 & OK \\
        \midrule
        \textbf{Total} & \textbf{1\,469} & \textbf{1\,469} & \textbf{OK} \\
        \bottomrule
      \end{tabular}
    \end{tablaAPA}%
  }{Elaboración propia con datos de \texttt{evidence/03\_pruebas\_y\_hardening.txt}.}
\end{table}

\subsection{Integridad referencial}

Se verificaron las 9 claves foráneas con consultas de huérfanos. La
Tabla~\ref{tab:fks} muestra las relaciones y el resultado. Todas devolvieron
0. La consulta genérica es \texttt{SELECT COUNT(*) FROM hija h LEFT JOIN padre
p ON h.fk = p.pk WHERE h.fk IS NOT NULL AND p.pk IS NULL}. Si el resultado es
distinto de cero, la clave foránea no se puede declarar o la carga tiene un
error.

\begin{table}[H]
  \centering
  \caption{Claves foráneas verificadas}
  \label{tab:fks}
  \contenidoConFuente{%
    \begin{tablaAPA}
      \begin{tabular}{@{}L{4.4cm}L{4.4cm}R{2.2cm}@{}}
        \toprule
        \textbf{Tabla hija} & \textbf{Tabla padre} & \textbf{Huérfanos} \\
        \midrule
        \texttt{cuentas\_ahorro.socio\_id} & \texttt{socios.id} & 0 \\
        \texttt{solicitudes\_credito.socio\_id} & \texttt{socios.id} & 0 \\
        \texttt{solicitudes\_credito.tipo\_credito\_id} & \texttt{tipos\_credito.id} & 0 \\
        \texttt{prestamos.solicitud\_id} & \texttt{solicitudes\_credito.id} & 0 \\
        \texttt{prestamos.socio\_id} & \texttt{socios.id} & 0 \\
        \texttt{cuotas\_amortizacion.prestamo\_id} & \texttt{prestamos.id} & 0 \\
        \texttt{transacciones.socio\_id} & \texttt{socios.id} & 0 \\
        \texttt{transacciones.cuenta\_ahorro\_id} & \texttt{cuentas\_ahorro.id} & 0 \\
        \texttt{transacciones.prestamo\_id} & \texttt{prestamos.id} & 0 \\
        \bottomrule
      \end{tabular}
    \end{tablaAPA}%
  }{Elaboración propia con datos de \texttt{evidence/03\_pruebas\_y\_hardening.txt}.}
\end{table}

\subsection{Rendimiento de consultas críticas}

Se midieron las dos vistas principales. La Tabla~\ref{tab:rendimiento}
distingue el tiempo de pared (\emph{wall}) del tiempo de servidor. El primero
incluye el costo de \texttt{docker compose exec}, la autenticación, la
transferencia y el formateo; el segundo es el \texttt{Execution Time} de
\texttt{EXPLAIN ANALYZE} y refleja el trabajo real del motor. Para
\texttt{vw\_cartera\_mora}, el servidor tardó 0.236 ms; para
\texttt{vw\_resumen\_saldos\_socio}, 0.317 ms.

\begin{table}[H]
  \centering
  \caption{Rendimiento de las vistas críticas}
  \label{tab:rendimiento}
  \contenidoConFuente{%
    \begin{tablaAPA}
      \begin{tabular}{@{}L{5.0cm}R{2.6cm}R{2.6cm}R{2.8cm}@{}}
        \toprule
        \textbf{Consulta} & \textbf{MySQL wall (ms)} & \textbf{PG wall (ms)} & \textbf{PG servidor (ms)} \\
        \midrule
        \texttt{vw\_cartera\_mora} & 124 & 143 & 0.236 \\
        \texttt{vw\_resumen\_saldos\_socio} & 108 & 137 & 0.317 \\
        \bottomrule
      \end{tabular}
    \end{tablaAPA}%
  }{Elaboración propia con datos de \texttt{evidence/03\_pruebas\_y\_hardening.txt}.}
\end{table}

El plan de \texttt{vw\_cartera\_mora} usa un \texttt{Bitmap Index Scan} sobre
\texttt{idx\_prestamo\_estado} para localizar los préstamos morosos, luego un
\texttt{Hash Join} con socios y cuotas, y un \texttt{GroupAggregate}. El plan
de \texttt{vw\_resumen\_saldos\_socio} usa \texttt{Merge Left Join} e
\texttt{Incremental Sort}. Para el volumen actual, los \texttt{Seq Scan} sobre
tablas pequeñas son razonables; a mayor escala, el optimizador podría cambiar
a índices si las estadísticas y la selectividad lo justifican. La lección es
que el rendimiento debe medirse con \texttt{EXPLAIN (ANALYZE, BUFFERS)} y no
suponerse por el tamaño de la tabla.

\section{Confidencialidad y escalabilidad}

\subsection{Clasificación de la información}

La Tabla~\ref{tab:clasificacion} clasifica los datos tratados. Los datos
personales y financieros son confidenciales o restringidos. La clasificación
guía las medidas: los datos restringidos solo se exponen a través de vistas
enmascaradas; los confidenciales se cifran en tránsito y en reposo; los
internos se protegen con roles; y los públicos pueden compartirse sin
restricción adicional.

\begin{table}[H]
  \centering
  \caption{Clasificación de la información}
  \label{tab:clasificacion}
  \contenidoConFuente{%
    \begin{tablaAPA}
      \begin{tabular}{@{}L{2.8cm}L{5.0cm}L{6.0cm}@{}}
        \toprule
        \textbf{Nivel} & \textbf{Ejemplos} & \textbf{Control principal} \\
        \midrule
        Restringido & CI/NIT, número de cuenta, teléfono, correo, dirección. & Enmascaramiento en \texttt{vw\_socios\_audit}; sin acceso directo para \texttt{auditor\_ro}. \\
        Confidencial & Saldos, montos de crédito, cuotas, transacciones. & TLS, roles, respaldo cifrado, auditoría. \\
        Interno & Catálogos de sucursales y tipos de crédito. & Roles y registros de acceso. \\
        Público & Documentación, scripts, informe. & Sin restricción adicional. \\
        \bottomrule
      \end{tabular}
    \end{tablaAPA}%
  }{Elaboración propia.}
\end{table}

\subsection{Controles de confidencialidad}

Los controles se organizan en cuatro momentos: en tránsito, en reposo, en uso
y en auditoría. En tránsito, TLS 1.2+ con \texttt{hostssl} y
\texttt{scram-sha-256}. En reposo, AES-256-CBC para respaldos y permisos
\texttt{600} para la clave privada. En uso, roles de mínimo privilegio y
vistas enmascaradas. En auditoría, registros de conexión, desconexión y DDL,
más la separación del rol \texttt{auditor\_ro}. La combinación es más fuerte
que cualquier control aislado: si un atacante obtiene una credencial de
aplicación, no puede leer la tabla de socios; si obtiene un respaldo, no puede
descifrarlo sin la frase de paso; si intenta conectarse sin TLS, el servidor
lo rechaza.

\subsection{Escalabilidad}

El diseño actual es adecuado para el volumen de la cooperativa, pero la
escalabilidad requiere un plan. La Tabla~\ref{tab:escalabilidad} propone
medidas por horizonte. A corto plazo, se recomienda añadir un \emph{pool} de
conexiones (PgBouncer), monitorear consultas lentas y ajustar
\texttt{autovacuum}. A mediano plazo, particionar \texttt{transacciones} por
fecha y crear índices parciales para los estados más consultados. A largo
plazo, evaluar réplicas de lectura para reportes y una arquitectura de alta
disponibilidad con \emph{failover}.

\begin{table}[H]
  \centering
  \caption{Medidas de escalabilidad por horizonte}
  \label{tab:escalabilidad}
  \contenidoConFuente{%
    \begin{tablaAPA}
      \begin{tabular}{@{}L{2.4cm}L{6.2cm}L{5.6cm}@{}}
        \toprule
        \textbf{Horizonte} & \textbf{Medida} & \textbf{Beneficio} \\
        \midrule
        Corto plazo & PgBouncer, monitoreo de consultas lentas, ajuste de \texttt{autovacuum}. & Menos conexiones, mejor uso de recursos. \\
        Mediano plazo & Particionamiento de \texttt{transacciones} por rango de fechas; índices parciales. & Consultas históricas más rápidas; mantenimiento por partición. \\
        Largo plazo & Réplicas de lectura para reportes; alta disponibilidad con \emph{failover}. & Escalabilidad horizontal y continuidad. \\
        Transversal & Capacitación, documentación y pruebas de carga periódicas. & Operación sostenible y decisiones basadas en datos. \\
        \bottomrule
      \end{tabular}
    \end{tablaAPA}%
  }{Elaboración propia.}
\end{table}

La escalabilidad no es solo agregar hardware: es evitar cuellos de botella
conocidos. Un \emph{pool} mal configurado puede ser peor que no tenerlo; un
índice innecesario ralentiza las escrituras; una réplica retrasada puede
servir datos obsoletos. Por eso las medidas se proponen por horizonte y con
métricas de validación, no como una lista de deseos.

% ============================================================================
%  CAPÍTULO III: DESARROLLO DEL TRABAJO (parte 3)
% ============================================================================

\section{Consideraciones del sistema operativo destino}

\subsection{Aislamiento con contenedores}

El entorno se ejecuta sobre un host Linux con Docker Engine y Docker Compose
v2. Cada servicio corre en su propio contenedor, con su propio sistema de
archivos, su propia pila de red y sus propios límites de recursos. Los
contenedores comparten el núcleo del host, pero no comparten el espacio de
usuarios ni los procesos. Esa separación reduce el impacto de una
vulnerabilidad en un servicio y facilita la reproducibilidad
\citep{docker2024}.

Los puertos se publican únicamente en \texttt{127.0.0.1}. Esa decisión evita
que MySQL o PostgreSQL queden accesibles desde la red local o desde Internet.
Si un integrante necesita conectarse con una herramienta gráfica, puede usar
un túnel SSH o una regla explícita; no se abre el puerto a \texttt{0.0.0.0}
por comodidad. La red interna \texttt{red\_cooperativa} permite que el
migrador resuelva \texttt{db-origen} y \texttt{db-destino} por nombre, sin
exponer los motores fuera del proyecto. El contenedor del migrador está
construido sobre \texttt{postgres:16} con \texttt{pgloader} 3.6.10,
\texttt{gettext-base} y \texttt{tzdata}; ya no requiere Python ni controladores
en el host.

\subsection{Permisos de archivos y secretos}

La clave privada TLS \texttt{server.key} tiene permisos \texttt{600} y
pertenece al usuario \texttt{postgres}. El certificado \texttt{server.crt}
tiene permisos \texttt{644}. El archivo \texttt{.env} contiene las
contraseñas de laboratorio, no se versiona y está incluido en
\texttt{.gitignore}. Los respaldos se guardan en \texttt{backups/}, también
excluido del control de versiones. En producción, las credenciales deben
gestionarse con un gestor de secretos y no con un archivo de texto.

\begin{table}[H]
  \centering
  \caption{Consideraciones del sistema operativo y del host}
  \label{tab:so}
  \contenidoConFuente{%
    \begin{tablaAPA}
      \begin{tabular}{@{}L{3.4cm}L{5.4cm}L{5.4cm}@{}}
        \toprule
        \textbf{Área} & \textbf{Configuración aplicada} & \textbf{Recomendación de producción} \\
        \midrule
        Aislamiento & Contenedores separados; red interna; puertos en \texttt{127.0.0.1}. & Mantener; añadir políticas de red y \emph{firewall}. \\
        Permisos & \texttt{server.key} en \texttt{600}; \texttt{.env} fuera de git. & Gestor de secretos; rotación; auditoría de accesos. \\
        Recursos & Sin límites explícitos en el laboratorio. & Definir \texttt{cpus}, \texttt{mem\_limit} y alertas. \\
        Registros & \texttt{docker logs}; \texttt{log\_line\_prefix}. & Rotación, recolección central y retención. \\
        Actualizaciones & Imágenes fijadas por versión mayor. & Escaneo de vulnerabilidades y parches programados. \\
        Zona horaria & \texttt{America/La\_Paz} en ambos motores. & Documentar y validar en reportes contables. \\
        Respaldos & Volúmenes nombrados; respaldo lógico cifrado. & Copia fuera del sitio y pruebas de restauración. \\
        \bottomrule
      \end{tabular}
    \end{tablaAPA}%
  }{Elaboración propia.}
\end{table}

\subsection{Zona horaria y consistencia temporal}

Ambos motores se configuraron con la zona \texttt{America/La\_Paz}. En MySQL,
la zona afecta la interpretación de \texttt{DATETIME} y las funciones de
fecha; en PostgreSQL, afecta la visualización de \texttt{TIMESTAMPTZ} y el
valor de \texttt{now()}. La migración asigna la zona en el ETL y el servidor
la conserva en las consultas. Para una entidad financiera, la consistencia
temporal es parte de la integridad: un pago no debe caer en el día anterior
o siguiente por un desajuste de zona.

\subsection{Actualizaciones y cadena de suministro}

Las imágenes se fijan por versión mayor (\texttt{mysql:5.7} y
\texttt{postgres:16}). En producción se recomienda fijar una versión exacta
(\emph{digest}) para evitar cambios inesperados, escanear las imágenes con una
herramienta de vulnerabilidades y actualizar con un proceso controlado. La
imagen de documentación (\texttt{texlive/texlive}) solo se usa en el perfil
\texttt{docs} y no forma parte del entorno operativo.

\section{Bitácora de problemas, decisiones y lecciones aprendidas}
\label{sec:bitacora}

\subsection{Problemas encontrados}

La Tabla~\ref{tab:bitacora} documenta los problemas reales que surgieron
durante la implementación. No se trata de una lista teórica: cada fila tiene
un síntoma, una causa, una solución y una evidencia. Documentar los errores es
parte del valor del proyecto, porque evita que el próximo equipo repita el
mismo camino.

\begingroup
\setstretch{1}\small\sloppy\emergencystretch=3em\relax
\begin{longtable}{@{}C{0.8cm}L{3.4cm}L{3.8cm}L{5.4cm}@{}}
  \caption{Bitácora de problemas y soluciones}\label{tab:bitacora}\\
  \toprule
  \textbf{\#} & \textbf{Síntoma} & \textbf{Causa} & \textbf{Solución aplicada} \\
  \midrule
  \endfirsthead
  \multicolumn{4}{c}{\textit{Continuación de la Tabla~\ref{tab:bitacora}}}\\
  \toprule
  \textbf{\#} & \textbf{Síntoma} & \textbf{Causa} & \textbf{Solución aplicada} \\
  \midrule
  \endhead
  \bottomrule
  \endlastfoot
  1 & El origen no tenía tablas al iniciar. & El volumen se creó antes de cargar el esquema; el entrypoint no se ejecutó. & Se recreó el volumen con \texttt{down -v} y se cargó el esquema en el primer arranque. \\
  2 & El seed inicial fallaba por cadenas sin comillas. & SQL generado a mano con valores de texto sin entrecomillar. & Se reemplazó por un generador Python determinista con \texttt{seed=42}. \\
  3 & \texttt{TRUNCATE ... RESTART IDENTITY} fallaba por permisos. & \texttt{admin\_dba} no era propietario de las secuencias de identidad. & Se eliminó \texttt{RESTART IDENTITY} y se usó \texttt{setval} al final. \\
  4 & PostgreSQL reiniciaba por certificado faltante. & La configuración apuntaba a una ruta no escribible o inexistente. & El entrypoint copia los certificados a \texttt{/var/lib/postgresql/ssl} y ajusta permisos. \\
  5 & El healthcheck generaba errores de \texttt{peer}. & Se ejecutaba como \texttt{root} con \texttt{-U postgres} por socket. & Se mantuvo el healthcheck; el contenedor queda \emph{healthy}; los errores son de autenticación \texttt{peer}, no de disponibilidad. \\
  6 & \texttt{auditor\_ro} podía ver datos sensibles. & La vista \texttt{vw\_cartera\_mora} exponía \texttt{numero\_identificacion}. & Se revocó el acceso a esa vista y se dejó solo \texttt{vw\_socios\_audit}. \\
  7 & El acceso remoto de \texttt{postgres} era posible. & \texttt{pg\_hba} incluía una regla \texttt{host} amplia. & Se eliminó la regla remota; solo \texttt{local peer}. \\
  8 & Las fechas perdían zona horaria. & \texttt{DATETIME} no almacena zona; el ETL no la asignaba. & Se usó \texttt{ZoneInfo("America/La\_Paz")} y \texttt{TIMESTAMPTZ}. \\
  9 & Los nuevos \texttt{INSERT} chocaban con IDs existentes. & Las secuencias no se reajustaron después de insertar IDs explícitos. & \texttt{setval} con \texttt{MAX(id)} en cada tabla. \\
  10 & El ETL no era idempotente. & Una segunda ejecución duplicaba filas. & \texttt{TRUNCATE} previo y snapshot consistente; verificación de conteos. \\
  11 & \texttt{postgres:16} es Debian 13 (trixie). & El paquete disponible es pgloader 3.6.10, no 3.6.9. & Se usó pgloader 3.6.10 (misma rama estable 3.6.x). \\
  12 & pgloader 3.6.10 verifica el certificado con \texttt{sslmode=require}. & \texttt{--no-ssl-cert-verification} no lo desactiva. & Se usó \texttt{SSL\_CERT\_FILE=/etc/postgresql/ssl/server.crt}. \\
  13 & pgloader no tiene opción \texttt{create no schema} para MySQL. & Crea un schema con el nombre de la base origen. & Se resolvió con \texttt{ALTER SCHEMA ... RENAME TO 'stg'}. \\
  14 & \texttt{ANALYZE} exige ser dueño en PostgreSQL 16. & \texttt{MAINTAIN} recién aparece en PG 17. & Se ejecuta como \texttt{postgres} por socket local. \\
  15 & pgloader escribe rechazos en \texttt{/tmp/pgloader/} por defecto. & Ruta por defecto no configurable sin parámetro. & Se usó \texttt{--root-dir /tmp/pgloader-work}. \\
  16 & Tags flotantes sin fijar por digest. & \texttt{mysql:5.7} y \texttt{postgres:16} pueden cambiar. & Mejora pendiente para producción: fijar por digest. \\
\end{longtable}
\fuenteAPA{Elaboración propia con base en la ejecución real y \texttt{evidence/}.}
\endgroup

\subsection{Decisiones de diseño}

La Tabla~\ref{tab:decisiones} resume las decisiones más importantes y su
justificación. Algunas fueron técnicamente debatibles; se documentan para que
el lector pueda evaluarlas y, si es necesario, cambiarlas con conocimiento de
causa.

\begin{table}[H]
  \centering
  \caption{Decisiones de diseño y su justificación}
  \label{tab:decisiones}
  \contenidoConFuente{%
    \begin{tablaAPA}
      \begin{tabular}{@{}L{3.8cm}L{4.4cm}L{6.0cm}@{}}
        \toprule
        \textbf{Decisión} & \textbf{Alternativa descartada} & \textbf{Justificación} \\
        \midrule
        ETL propio en Python & pgloader u otra herramienta genérica & Control explícito de tipos, zona horaria, identidades y verificación; el volumen es pequeño. \\
        Conservar nombres de origen & Renombrar columnas a un estándar nuevo & Reduce el impacto en la aplicación y facilita la comparación. \\
        \texttt{TIMESTAMPTZ} en \texttt{America/La\_Paz} & \texttt{TIMESTAMP} sin zona & Evita ambigüedad en pagos y reportes contables. \\
        \texttt{admin\_dba} sin \texttt{SUPERUSER} & Superusuario para todo & Limita el daño de un error o compromiso. \\
        Sin \texttt{RESTART IDENTITY} & \texttt{TRUNCATE ... RESTART IDENTITY} & Evita el error de propiedad de secuencias; \texttt{setval} es explícito. \\
        Puertos en \texttt{127.0.0.1} & Publicar en \texttt{0.0.0.0} & Reduce la exposición; el acceso remoto se hace por túnel. \\
        Vista enmascarada & Copia anonimizada completa & Permite auditoría sin duplicar datos ni romper referencias. \\
        \bottomrule
      \end{tabular}
    \end{tablaAPA}%
  }{Elaboración propia.}
\end{table}

\subsection{Lecciones aprendidas}

La primera lección es que la migración empieza por el diagnóstico, no por la
copia. Sin inventario, sin conteos y sin matriz de incompatibilidades, el ETL
es una caja negra. La segunda es que la seguridad debe diseñarse desde el
principio: si se deja para el final, se convierte en una capa de parches. La
tercera es que las pruebas deben ser automatizadas y comparar valores
esperados, no solo imprimir resultados. La cuarta es que la documentación de
los errores es tan valiosa como la de los aciertos: el problema de
\texttt{TRUNCATE ... RESTART IDENTITY} habría costado horas a otro equipo. La
quinta es que la reproducibilidad no es un lujo: es la única forma de afirmar
que el resultado es confiable.

\section{Lista de cotejo de seguridad}

La lista de cotejo de seguridad aplicada al destino se presenta en el
Anexo~C, sección C.2 (Tabla~\ref{tab:checklist}). Cada control tiene un estado
y una evidencia: los marcados como ``Cumple'' se verificaron en el entorno y
los marcados como ``Recomendado'' son mejoras para producción fuera del
alcance del laboratorio.

\section{Trazabilidad de evidencias}

La Tabla~\ref{tab:trazabilidad} relaciona cada punto exigido con la sección
del informe, el archivo de evidencia y el resultado. Esta tabla es la prueba
de que el informe no es una narración independiente de los datos: cada
afirmación tiene un origen verificable.

\begin{table}[H]
  \centering
  \caption{Trazabilidad de los puntos exigidos}
  \label{tab:trazabilidad}
  \contenidoConFuente{%
    \begin{tablaAPA}
      \begin{tabular}{@{}C{0.7cm}L{4.6cm}L{3.8cm}L{4.6cm}@{}}
        \toprule
        \textbf{\#} & \textbf{Punto exigido} & \textbf{Evidencia} & \textbf{Resultado} \\
        \midrule
        1 & Portada, resumen, objetivos y alcance. & Este documento & Secciones 1.1--1.7. \\
        2 & Cliente, problema, solución y ODS. & Este documento & Secciones 1.2--1.5. \\
        3 & Diagnóstico del origen. & \texttt{evidence/01} & 8 tablas, 2 vistas, 0.28 MB. \\
        4 & Hardening y verificación. & \texttt{evidence/03} & 75 PASS, 0 FAIL. \\
        5 & Plan de migración y contingencia. & Este documento & Sección 3.1. \\
        6 & Ejecución y salidas reales. & \texttt{evidence/02} & 1\,469 filas, 0 diferencias. \\
        7 & Pruebas en destino. & \texttt{evidence/03} & Conteos, FK y rendimiento. \\
        8 & Confidencialidad y escalabilidad. & Este documento & Sección 3.8. \\
        9 & Consideraciones del SO. & Este documento & Sección 3.9. \\
        10 & Bitácora y lecciones. & Este documento & Sección 3.10. \\
        11 & Conclusiones, recomendaciones y bibliografía. & Este documento & Capítulo 5 y referencias. \\
        12 & Anexos y lista de cotejo. & Anexos A--C & Scripts completos. \\
        13 & Guion de demostración. & \texttt{GUION\_DEMO\_EN\_VIVO.md} & 15--20 min. \\
        \bottomrule
      \end{tabular}
    \end{tablaAPA}%
  }{Elaboración propia.}
\end{table}

\section{Reproducibilidad y transferencia}

\subsection{Procedimiento de ejecución}

El procedimiento completo se ejecuta en cuatro pasos. Primero, se copia el
archivo \texttt{.env.example} a \texttt{.env} y se ajustan las credenciales.
Segundo, se levanta el entorno con \texttt{docker compose up -d --build}.
Tercero, se ejecutan los scripts en orden. Cuarto, se revisan las evidencias y
se compila el informe. La secuencia exacta se muestra a continuación.

\begin{lstlisting}[style=terminal,caption={Secuencia de reproducción del proyecto},label={lst:reproduccion}]
cp .env.example .env
docker compose up -d --build
./scripts/01_audit_origin.sh
./scripts/02_migrate_and_backup.sh
./scripts/03_verify_all.sh
docker compose --profile docs run --rm docs make
\end{lstlisting}

\subsection{Resultados esperados}

Si el entorno se levanta desde cero, los resultados esperados son: 8 tablas y
2 vistas en el origen; 55 socios, 66 cuentas, 50 solicitudes, 32 préstamos,
516 cuotas y 740 transacciones; 1\,469 filas migradas; 75 pruebas en
\texttt{PASS}; y un PDF en \texttt{docs/main.pdf}. Cualquier diferencia debe
investigarse antes de aceptar la migración. La semilla fija del generador
garantiza que los datos sean los mismos en cada ejecución.

\subsection{Transferencia}

El proyecto se entrega con la siguiente estructura: \texttt{src/} (esquemas,
seed, ETL y configuración), \texttt{scripts/} (automatización),
\texttt{evidence/} (salidas reales), \texttt{backups/} (respaldos cifrados),
\texttt{docs/} (informe, anexos y guion) y \texttt{docker-compose.yml}. Las
credenciales de laboratorio están en \texttt{.env.example}; las reales no se
versionan. El lector puede reproducir el entorno en cualquier host Linux con
Docker y Docker Compose, sin instalar MySQL ni PostgreSQL en el sistema.


% --- CONCLUSIONES ------------------------------------------------------------
\newpage
\encabezadoCapitulo{CONCLUSIONES}{}
\addcontentsline{toc}{part}{CONCLUSIONES}
\begingroup
\setcounter{capitulo}{5}
\setcounter{section}{0}
\renewcommand{\thesection}{5.\arabic{section}}
\renewcommand{\theHsection}{5.\arabic{section}}
\renewcommand{\thesubsection}{5.\arabic{section}.\arabic{subsection}}
\renewcommand{\theHsubsection}{5.\arabic{section}.\arabic{subsection}}
\renewcommand{\thesubsubsection}{5.\arabic{section}.\arabic{subsection}.\arabic{subsubsection}}
\renewcommand{\theHsubsubsection}{5.\arabic{section}.\arabic{subsection}.\arabic{subsubsection}}
% ============================================================================
%  CONCLUSIONES
% ============================================================================

\section{Conclusiones}

El proyecto de migración heterogénea de la base de datos de la Cooperativa de
Ahorro y Crédito Horizonte PyME S.R.L. desde MySQL 5.7 hacia PostgreSQL 16 se
ha ejecutado de forma íntegra, segura y verificable, cumpliendo con los ocho
objetivos específicos definidos en el Capítulo~II.

\begin{enumerate}
  \item \textbf{Diagnóstico del sistema de origen:} se levantó un inventario
        completo del motor MySQL 5.7.44: 8 tablas, 2 vistas, 19 entradas de
        índice, 0.28 MB de tamaño, un único usuario de aplicación
        (\texttt{app\_legacy\_user} con \texttt{GRANT ALL}), ausencia de claves
        foráneas y tráfico sin cifrado.
  \item \textbf{Incompatibilidades de tipos y semántica:} se documentó la matriz
        de conversión (\texttt{TINYINT(1)} $\rightarrow$ \texttt{BOOLEAN},
        \texttt{ENUM} $\rightarrow$ \texttt{VARCHAR} + \texttt{CHECK},
        \texttt{DATETIME} $\rightarrow$ \texttt{TIMESTAMPTZ} en
        \texttt{America/La\_Paz}, \texttt{INT UNSIGNED AUTO\_INCREMENT}
        $\rightarrow$ \texttt{BIGINT GENERATED ALWAYS AS IDENTITY}).
  \item \textbf{Diseño del esquema destino:} se implementaron tipos nativos de
        PostgreSQL 16, restricciones \texttt{CHECK}, 9 claves foráneas, 10
        índices y 3 vistas (incluyendo la vista enmascarada
        \texttt{vw\_socios\_audit}).
  \item \textbf{ETL determinista:} se reemplazó el ETL artesanal en Python por
        una arquitectura basada en \textbf{pgloader} con schema de staging
        (\texttt{stg}), precheck de validación, carga en una sola transacción con
        \texttt{OVERRIDING SYSTEM VALUE}, postcheck bidireccional
        (\texttt{to\_jsonb} + \texttt{EXCEPT ALL}), reajuste de secuencias con
        \texttt{setval} y \texttt{ANALYZE}. El proceso es idempotente y con
        rollback transaccional garantizado.
  \item \textbf{Endurecimiento de PostgreSQL 16:} TLS 1.2 obligatorio con
        \texttt{hostssl}, autenticación \texttt{scram-sha-256}, revocación de
        \texttt{PUBLIC}, 7 roles con mínimo privilegio
        (\texttt{admin\_dba}, \texttt{app\_core\_rw}, \texttt{auditor\_ro},
        \texttt{migration\_loader} y 4 roles de administradores del equipo),
        registro de conexiones, desconexiones y DDL.
  \item \textbf{Enmascaramiento de datos sensibles:} la vista
        \texttt{vw\_socios\_audit} expone CI/NIT, teléfono, correo y números de
        cuenta enmascarados; \texttt{auditor\_ro} no tiene acceso directo a
        tablas base sensibles.
  \item \textbf{Respaldo cifrado y restauración:} \texttt{pg\_dump -Fc} cifrado
        con AES-256-CBC y PBKDF2 (200.000 iteraciones), verificación SHA-256 y
        prueba de restauración en base temporal con los 55 socios recuperados.
  \item \textbf{Automatización y evidencias:} scripts reproducibles
        (\texttt{01\_audit\_origin.sh}, \texttt{02\_migrate\_and\_backup.sh},
        \texttt{03\_verify\_all.sh}) con salidas reales en \texttt{evidence/}.
\end{enumerate}

El proyecto cumple con los objetivos planteados: la base de datos de la
cooperativa ahora opera sobre PostgreSQL 16, un motor moderno y soportado, con
seguridad por capas, integridad referencial, datos sensibles enmascarados,
respaldo cifrado verificable y un proceso de migración reproducible y
automatizado.

\section{Recomendaciones}

A partir de la ejecución y de las pruebas realizadas se recomienda, para una
adopción en producción:

\begin{enumerate}
  \item \textbf{Gestión centralizada de secretos:} sustituir \texttt{.env} por un
        gestor (Vault, AWS Secrets Manager o equivalente), con rotación
        periódica de la frase de paso de los respaldos y separación entre la
        custodia de la clave y la del archivo cifrado.
  \item \textbf{Respaldo fuera del sitio:} aplicar la regla 3-2-1 y automatizar
        la verificación SHA-256 y la prueba de restauración periódica, no solo
        la generación del respaldo.
  \item \textbf{Reducción de privilegios en el origen:} sustituir el
        \texttt{GRANT ALL} de \texttt{app\_legacy\_user} por permisos mínimos de
        lectura antes de cualquier nueva corrida sobre datos reales.
  \item \textbf{Reproducibilidad de imágenes:} fijar las imágenes
        \texttt{mysql:5.7} y \texttt{postgres:16} por digest para evitar cambios
        silenciosos entre despliegues.
  \item \textbf{Monitoreo y capacidad:} instrumentar conexiones, consultas
        lentas y crecimiento de los datos; evaluar PgBouncer y, cuando el
        volumen lo justifique, replicación con CDC para reducir la ventana de
        corte.
  \item \textbf{Operación formal:} documentar un \emph{runbook} de reversión y
        ejecutar ensayos periódicos de recuperación con RTO y RPO medidos.
\end{enumerate}

\endgroup

% --- Referencias -------------------------------------------------------------
\newpage
\encabezadoCapitulo{REFERENCIAS}{}
\addcontentsline{toc}{part}{REFERENCIAS}
\bibliographystyle{apacite}
\nocite{*}
\bibliography{formato/referencias}

% --- Anexos ------------------------------------------------------------------
\newpage
\encabezadoCapitulo{ANEXOS}{}
\pdfbookmark[0]{ANEXOS}{anexos}%
\addcontentsline{toc}{part}{ANEXOS}
\appendix
% ============================================================================
%  ANEXO A: ESQUEMAS SQL DEL ORIGEN Y DEL DESTINO
% ============================================================================

\section{Esquemas SQL del origen y del destino (Anexo A)}

Este anexo reproduce los archivos SQL que definen el sistema de origen y el
sistema de destino. Se incluyen completos para que el lector pueda recrear el
entorno sin depender de fragmentos. El esquema de origen corresponde a MySQL
5.7 y contiene las incompatibilidades deliberadas descritas en el
Capítulo~III; el esquema de destino corresponde a PostgreSQL 16 e incorpora
tipos nativos, restricciones \texttt{CHECK}, claves foráneas e identidades.

\subsection{Esquema del origen (MySQL 5.7)}

El archivo \texttt{src/origin/01\_schema\_origin.sql} crea la base
\texttt{coop\_horizonte\_legacy}, las 8 tablas, las 2 vistas, los índices y el
usuario \texttt{app\_legacy\_user} con \texttt{GRANT ALL}. Obsérvese el uso de
\texttt{TINYINT(1)}, \texttt{ENUM}, \texttt{DATETIME}, \texttt{INT UNSIGNED
AUTO\_INCREMENT} y \texttt{utf8\_general\_ci}.

\lstinputlisting[language=SQL,caption={Esquema completo del origen MySQL 5.7},label={lst:anexo-origen}]{anexos/codigo/01_schema_origin.sql}

\subsection{Esquema del destino (PostgreSQL 16)}

El archivo \texttt{src/target/01\_schema\_target.sql} crea las 8 tablas con
identidades \texttt{BIGINT}, booleanos, marcas de tiempo con zona horaria,
restricciones \texttt{CHECK}, 9 claves foráneas, 10 índices y las 3 vistas. La
zona horaria se fija en \texttt{America/La\_Paz}.

\lstinputlisting[language=SQL,caption={Esquema completo del destino PostgreSQL 16},label={lst:anexo-destino}]{anexos/codigo/01_schema_target.sql}

\subsection{Roles y privilegios}

El archivo \texttt{src/target/02\_hardening\_roles.sql} revoca los privilegios
de \texttt{PUBLIC}, crea los roles \texttt{admin\_dba}, \texttt{app\_core\_rw}
y \texttt{auditor\_ro}, y aplica la matriz de mínimo privilegio. También crea
la vista enmascarada \texttt{vw\_socios\_audit}.

\lstinputlisting[language=SQL,caption={Endurecimiento de roles y privilegios},label={lst:anexo-roles}]{anexos/codigo/02_hardening_roles.sql}

% ============================================================================
%  ANEXO B: SCRIPTS DE AUTOMATIZACIÓN, ETL Y ORQUESTACIÓN
% ============================================================================

\section{Scripts de automatización, ETL y orquestación (Anexo B)}

Este anexo reproduce los scripts y archivos de configuración que hacen
reproducible el proyecto. Todos los scripts cargan las variables desde
\texttt{.env}, esperan a que los servicios estén \emph{healthy} y escriben su
salida en \texttt{evidence/}. El ETL se ejecuta dentro del contenedor
\texttt{migrador}; no requiere instalar Python ni controladores en el host.

\subsection{Orquestación con Docker Compose}

El archivo \texttt{docker-compose.yml} define los servicios
\texttt{db-origen}, \texttt{db-destino} y \texttt{migrador}, además del perfil
\texttt{docs} para compilar el informe. Los puertos se publican solo en
\texttt{127.0.0.1}.

\lstinputlisting[language=bash,caption={Orquestación completa con Docker Compose},label={lst:anexo-compose}]{anexos/codigo/docker-compose.yml.txt}

\subsection{Configuración de PostgreSQL}

Los archivos \texttt{postgresql.conf} y \texttt{pg\_hba.conf} establecen TLS
obligatorio, autenticación \texttt{scram-sha-256}, auditoría de conexiones y
DDL, zona horaria y bloqueo del acceso remoto de \texttt{postgres}.

\lstinputlisting[language=bash,caption={Configuración del servidor PostgreSQL 16},label={lst:anexo-postgresql-conf}]{anexos/codigo/postgresql.conf}

\lstinputlisting[language=bash,caption={Reglas efectivas de \texttt{pg\_hba.conf}},label={lst:anexo-pg-hba}]{anexos/codigo/pg_hba.conf}

\subsection{Biblioteca común de los scripts}

El archivo \texttt{scripts/lib.sh} centraliza la carga del entorno, la espera
de servicios, la ejecución de consultas y los conteos. Evita duplicar lógica
entre los tres scripts principales.

\lstinputlisting[language=bash,caption={Biblioteca común \texttt{lib.sh}},label={lst:anexo-lib}]{anexos/codigo/lib.sh}

\subsection{Auditoría del origen}

El script \texttt{01\_audit\_origin.sh} genera el inventario del origen:
versión, tamaño, tablas, vistas, índices, usuarios, privilegios e
incompatibilidades. Su salida se guarda en
\texttt{evidence/01\_inventario\_origen.txt}.

\lstinputlisting[language=bash,caption={Script de auditoría del origen},label={lst:anexo-audit}]{anexos/codigo/01_audit_origin.sh}

\subsection{Migración y respaldo cifrado}

El script \texttt{02\_migrate\_and\_backup.sh} respalda el origen con
\texttt{mysqldump}, ejecuta el ETL, respalda el destino con \texttt{pg\_dump
-Fc}, cifra el resultado con AES-256-CBC y prueba la restauración en una base
temporal.

\lstinputlisting[language=bash,caption={Script de migración y respaldo cifrado},label={lst:anexo-migrate}]{anexos/codigo/02_migrate_and_backup.sh}

\subsection{Verificación integral}

El script \texttt{03\_verify\_all.sh} ejecuta las 75 comprobaciones de
hardening, roles, enmascaramiento, conteos, integridad referencial, staging y
secuencias. Es el archivo que produce la evidencia
\texttt{03\_pruebas\_y\_hardening.txt}.

\lstinputlisting[language=bash,caption={Script de verificación integral},label={lst:anexo-verify}]{anexos/codigo/03_verify_all.sh}

\subsection{Plantilla pgloader}

El archivo \texttt{src/migration/pgloader/mysql\_to\_stage.load} es la plantilla
DSL de pgloader que carga las 8 tablas desde MySQL directamente al schema de
staging \texttt{stg} (sin índices, sin claves foráneas, sin secuencias). Las
credenciales se inyectan mediante variables de entorno con
\texttt{envsubst}.

\lstinputlisting[language=SQL,caption={Plantilla pgloader: MySQL a staging},label={lst:anexo-pgloader-load}]{anexos/codigo/mysql_to_stage.load}

\subsection{Renderizado de la plantilla pgloader}

El script \texttt{src/migration/render\_load.sh} renderiza la plantilla
pgloader sustituyendo solo las variables conocidas, evitando que
\texttt{envsubst} reemplace cualquier \$\{...\} que aparezca en el DSL.

\lstinputlisting[language=bash,caption={Renderizado de la plantilla pgloader},label={lst:anexo-render-load}]{anexos/codigo/render_load.sh}

\subsection{SQL de staging y migración}

Los archivos SQL que se ejecutan en secuencia durante la migración:
\emph{prepare} (crea \texttt{stg}), \emph{precheck} (valida antes de tocar
\texttt{public}), \emph{load} (carga con \texttt{OVERRIDING SYSTEM VALUE}),
\emph{postcheck} (compara bidireccional con \texttt{to\_jsonb} + \texttt{EXCEPT
ALL}), \emph{sequences} (\texttt{setval}), \emph{analyze} y \emph{cleanup}
(elimina \texttt{stg}).

\lstinputlisting[language=SQL,caption={00: Preparación del schema de staging},label={lst:anexo-sql-00}]{anexos/codigo/00_prepare_staging.sql}

\lstinputlisting[language=SQL,caption={10: Precheck de validación de staging},label={lst:anexo-sql-10}]{anexos/codigo/10_precheck.sql}

\lstinputlisting[language=SQL,caption={20: Carga stg → public con OVERRIDING SYSTEM VALUE},label={lst:anexo-sql-20}]{anexos/codigo/20_load.sql}

\lstinputlisting[language=SQL,caption={30: Postcheck bidireccional to\_jsonb + EXCEPT ALL},label={lst:anexo-sql-30}]{anexos/codigo/30_postcheck.sql}

\lstinputlisting[language=SQL,caption={35: Reajuste de secuencias con setval},label={lst:anexo-sql-35}]{anexos/codigo/35_sequences.sql}

\lstinputlisting[language=SQL,caption={40: ANALYZE de estadísticas del planificador},label={lst:anexo-sql-40}]{anexos/codigo/40_analyze.sql}

\lstinputlisting[language=SQL,caption={50: Limpieza y eliminación del schema staging},label={lst:anexo-sql-50}]{anexos/codigo/50_cleanup.sql}

\subsection{ETL en Python (archivado como legado)}

El archivo \texttt{src/migration/legacy/migrate.py} era el ETL artesanal en
Python (\texttt{pymysql} + \texttt{psycopg2}) que se usó inicialmente. Fue
archivado como legado y ya no se ejecuta; la migración actual usa pgloader +
SQL declarativo (secciones anteriores).

\lstinputlisting[language=Python,caption={ETL completo en Python (legado, archivado)},label={lst:anexo-etl}]{anexos/codigo/migrate.py}

\lstinputlisting[language=bash,caption={Imagen del contenedor migrador},label={lst:anexo-dockerfile}]{anexos/codigo/Dockerfile}

% ============================================================================
%  ANEXO C: CHECKLIST DE ACEPTACIÓN Y EVIDENCIAS
% ============================================================================

\section{Checklist de aceptación y evidencias (Anexo C)}

Este anexo reúne la lista de cotejo de aceptación y las evidencias reales
generadas por los scripts. Las evidencias se presentan como capturas de la
salida real de los scripts; los archivos de texto originales se conservan en
\texttt{evidence/} y cualquier lector puede contrastarlos con la salida de los
scripts descritos en el Anexo~B.

\subsection{Checklist de aceptación}

La Tabla~\ref{tab:checklist-aceptacion} enumera los criterios que debían
cumplirse para aceptar la migración. Todos los criterios obligatorios se
cumplieron; los recomendados quedaron documentados como trabajo futuro.

\begin{table}[H]
  \centering
  \caption{Checklist de aceptación de la migración}
  \label{tab:checklist-aceptacion}
  \contenidoConFuente{%
    \begin{tablaAPA}
      \begin{tabular}{@{}C{0.7cm}L{7.0cm}C{2.2cm}L{4.4cm}@{}}
        \toprule
        \textbf{\#} & \textbf{Criterio} & \textbf{Estado} & \textbf{Evidencia} \\
        \midrule
        1 & Entorno levantado con un solo comando. & Cumple & \texttt{docker compose up}. \\
        2 & 8 tablas y 2 vistas creadas en el origen. & Cumple & Anexo~C, Figs.~\ref{fig:anexo-inventario-1}--\ref{fig:anexo-inventario-2}. \\
        3 & 1\,469 filas migradas sin pérdida. & Cumple & Anexo~C, Figs.~\ref{fig:anexo-migracion-1}--\ref{fig:anexo-migracion-2}. \\
        4 & Conteos idénticos en las 8 tablas. & Cumple & Anexo~C, Figs.~\ref{fig:anexo-pruebas-1}--\ref{fig:anexo-pruebas-3}. \\
        5 & 9 claves foráneas sin huérfanos. & Cumple & Anexo~C, Figs.~\ref{fig:anexo-pruebas-1}--\ref{fig:anexo-pruebas-3}. \\
        6 & TLS obligatorio y conexiones sin SSL rechazadas. & Cumple & Anexo~C, Figs.~\ref{fig:anexo-pruebas-1}--\ref{fig:anexo-pruebas-3}. \\
        7 & Autenticación \texttt{scram-sha-256}. & Cumple & Anexo~C, Figs.~\ref{fig:anexo-pruebas-1}--\ref{fig:anexo-pruebas-3}. \\
        8 & Tres roles de mínimo privilegio. & Cumple & Anexo~C, Figs.~\ref{fig:anexo-pruebas-1}--\ref{fig:anexo-pruebas-3}. \\
        9 & \texttt{auditor\_ro} sin acceso a tablas sensibles. & Cumple & Anexo~C, Figs.~\ref{fig:anexo-pruebas-1}--\ref{fig:anexo-pruebas-3}. \\
        10 & Enmascaramiento de CI/NIT y cuentas. & Cumple & Anexo~C, Figs.~\ref{fig:anexo-pruebas-1}--\ref{fig:anexo-pruebas-3}. \\
        11 & Respaldo cifrado con AES-256-CBC. & Cumple & \texttt{backups/}. \\
        12 & Restauración verificada con 55 socios. & Cumple & Anexo~C, Figs.~\ref{fig:anexo-migracion-1}--\ref{fig:anexo-migracion-2}. \\
        13 & 75 pruebas en \texttt{PASS} y 0 en \texttt{FAIL}. & Cumple & Anexo~C, Figs.~\ref{fig:anexo-pruebas-1}--\ref{fig:anexo-pruebas-3}. \\
        14 & Informe compilado y guion de demo. & Cumple & \texttt{docs/}. \\
        15 & Copia de respaldo fuera del sitio. & Recomendado & Trabajo futuro. \\
        16 & Monitoreo, alertas y PgBouncer. & Recomendado & Trabajo futuro. \\
        \bottomrule
      \end{tabular}
    \end{tablaAPA}%
  }{Elaboración propia con base en la ejecución real del proyecto.}
\end{table}

\subsection{Evidencia 1: inventario del origen}

\capturaAPA{fig:anexo-inventario-1}{Inventario del origen: versión, tablas y conteos}%
  {captura-anexo-inventario-1.png}%
  {Salida de \texttt{scripts/01\_audit\_origin.sh} con la versión de MySQL, las tablas y sus conteos.}

\capturaAPA{fig:anexo-inventario-2}{Inventario del origen: índices, usuarios y privilegios}%
  {captura-anexo-inventario-2.png}%
  {Salida del inventario con los índices, el usuario \texttt{app\_legacy\_user} y sus \texttt{GRANT ALL}.}

\subsection{Evidencia 2: ejecución de la migración}

\capturaAPA{fig:anexo-migracion-1}{Ejecución de la migración: carga y conteos}%
  {captura-anexo-migracion-1.png}%
  {Salida del ETL con la carga de las 8 tablas y la verificación de conteos origen--destino.}

\capturaAPA{fig:anexo-migracion-2}{Respaldo cifrado y restauración}%
  {captura-anexo-migracion-2.png}%
  {Respaldo cifrado con AES-256-CBC y restauración verificada con 55 socios.}

\subsection{Evidencia 3: pruebas y hardening}

\capturaAPA{fig:anexo-pruebas-1}{Pruebas de SSL, \texttt{pg\_hba} y roles}%
  {captura-anexo-pruebas-1.png}%
  {Salida de \texttt{scripts/03\_verify\_all.sh} con el rechazo de conexiones sin SSL y la matriz de roles.}

\capturaAPA{fig:anexo-pruebas-2}{Enmascaramiento, conteos e integridad referencial}%
  {captura-anexo-pruebas-2.png}%
  {Salida de las pruebas de enmascaramiento, conteos origen--destino e integridad de las 9 claves foráneas.}

\capturaAPA{fig:anexo-pruebas-3}{Resumen de la verificación}%
  {captura-anexo-pruebas-3.png}%
  {Resumen final de \texttt{scripts/03\_verify\_all.sh}: \texttt{PASS: 75} y \texttt{FAIL: 0}.}


\end{document}
